\documentclass[10pt,a4paper]{article}
\usepackage[margin=16mm,headsep=5mm,footskip=8mm]{geometry}
\usepackage{amsmath,amssymb,booktabs,tabularx,array,graphicx,xcolor,enumitem,fancyhdr,hyperref,float,caption}
\usepackage[most]{tcolorbox}
\definecolor{navy}{HTML}{17324D}\definecolor{blue}{HTML}{2D6A8A}\definecolor{pale}{HTML}{EEF5F8}\definecolor{warn}{HTML}{8A4B08}
\hypersetup{colorlinks=true,linkcolor=blue,urlcolor=blue}
\setlength{\parindent}{0pt}\setlength{\parskip}{4pt}
\setlist{nosep,leftmargin=*}
\pagestyle{fancy}\fancyhf{}\lhead{ENME601 Thermodynamics}\rhead{Tutorial Booklet v2}\cfoot{\thepage}
\newtcolorbox{method}[1][]{colback=pale,colframe=blue,boxrule=.5pt,arc=1mm,title=#1,fonttitle=\bfseries}
\newtcolorbox{caution}[1][]{colback=yellow!8,colframe=warn,boxrule=.5pt,arc=1mm,title=#1,fonttitle=\bfseries}
\newcommand{\questionfigure}[2]{\begin{figure}[H]\centering\includegraphics[width=.88\textwidth,height=.43\textheight,keepaspectratio]{question-images/#1}\caption*{\small #2}\end{figure}}
\newcommand{\weekstart}[2]{\clearpage\phantomsection\addcontentsline{toc}{section}{#1 --- #2}}
\begin{document}
\begin{titlepage}\centering\vspace*{24mm}
{\Huge\bfseries Thermodynamics Tutorial Booklet v2\par}\vspace{5mm}
{\Large Weeks 1--6 --- Rewritten Questions and Expanded Solutions\par}\vspace{11mm}
{\large Every recoverable question is written out with the answer beside it, including system type, boundary, assumptions, property model, equations, substitutions, source caveats and necessary figures.\par}
\vfill
\begin{method}[How to use this booklet]
Read the rewritten question first. Before looking at the arithmetic, identify the boundary and write the unsimplified mass and energy balances. The ``vague/source note'' records ambiguities rather than silently repairing them. Numerical properties are tied to the current local tutorial tables and official solution pages; use the separate 2026 property tables during the exam for any fresh calculation.
\end{method}
\vfill 20 August 2026
\end{titlepage}
\tableofcontents\newpage
\section{A repeatable answer structure}
\begin{enumerate}
\item \textbf{System and boundary:} closed system/control mass, open control volume, or isolated system.
\item \textbf{Given and find:} convert litres, centimetres, gauge pressure and Celsius before substitution.
\item \textbf{Assumptions:} steady/transient, adiabatic, rigid, ideal gas, incompressible, negligible kinetic/potential change, quasi-equilibrium.
\item \textbf{State model:} classify the phase and select a saturation, superheated, compressed-liquid or ideal-gas table.
\item \textbf{Governing laws:} mass balance first for open systems, then the full energy balance.
\item \textbf{Working:} show table values, interpolation, units and sign convention.
\item \textbf{Check:} conservation, bounds, sign, physical direction and plausible magnitude.
\end{enumerate}
\begin{caution}[Notation]
This booklet uses $Q>0$ into the system and $W>0$ by the system unless a question asks directly for work input. A boxed positive compressor-work value is therefore reported as a required input magnitude.
\end{caution}
\weekstart{Week 1}{Basic properties, pressure and hydrostatics}
\section*{Week 1}

\subsection*{Question 1}
\textbf{Rewritten question.} The mass of air in a room of dimensions $3\,\text{m} \times 4\,\text{m} \times 25\,\text{m}$ is $350\,\text{kg}$. Determine the density, specific volume, and specific weight.

\textbf{Source wording caveat.} None.

\textbf{System type and boundary.} Closed quantity of air occupying the room volume; boundary is the room enclosure.

\textbf{Assumptions.} Air is treated as a continuum; gravity is uniform at $g=9.81\,\text{m/s}^2$.

\textbf{Given / Find.}
Given: $m=350\,\text{kg}$, $V=3\times 4\times 25=300\,\text{m}^3$.
Find: $\rho$, $v$, $\gamma$.

\textbf{Governing equations.}
\[
\rho=\frac{m}{V},\qquad v=\frac{1}{\rho},\qquad \gamma=\rho g
\]

\textbf{Sequential substitution with units.}
\[
\rho=\frac{350\,\text{kg}}{300\,\text{m}^3}=1.1667\,\text{kg/m}^3
\]
\[
v=\frac{1}{1.1667\,\text{kg/m}^3}=0.8571\,\text{m}^3/\text{kg}
\]
\[
\gamma=(1.1667\,\text{kg/m}^3)(9.81\,\text{m/s}^2)=11.44\,\text{N/m}^3
\]

\textbf{Final answer.}
\[
\boxed{\rho=1.167\,\text{kg/m}^3,\quad v=0.857\,\text{m}^3/\text{kg},\quad \gamma=11.44\,\text{N/m}^3}
\]

\textbf{Interpretation / common trap.} Density is mass per volume, not weight per volume; specific weight is $\rho g$, not $g/\rho$.

\textbf{Source page.} Page 5.

\subsection*{Question 2}
\textbf{Rewritten question.} Complete the missing properties in the table below if $g=9.81\,\text{m/s}^2$ and $V=104\,\text{L}$.

\textbf{Source wording caveat.} The original prompt is a fill-in table; the provided solution deck gives the completed entries but not a fully typeset blank table in text. The recoverable task is the table completion itself.

\textbf{System type and boundary.} Property-conversion exercise for a fixed amount of matter; no process boundary analysis is required.

\textbf{Assumptions.} Use $1\,\text{L}=10^{-3}\,\text{m}^3$ and $W=mg$.

\textbf{Given / Find.}
Given: $V=104\,\text{L}=0.104\,\text{m}^3$, $g=9.81\,\text{m/s}^2$.
Find: the consistent set of $(v,\rho,\gamma,m,W)$ for the five cases shown in the source solution.

\textbf{Governing equations.}
\[
\rho=\frac{1}{v},\qquad m=\rho V,\qquad W=mg,\qquad \gamma=\rho g
\]

\textbf{Sequential substitution with units.}
For each row, the solution deck already provides the computed values. Reconstructing the completed table:
\[
\begin{array}{c|ccccc}
\text{case} & v\,(\text{m}^3/\text{kg}) & \rho\,(\text{kg/m}^3) & \gamma\,(\text{N/m}^3) & m\,(\text{kg}) & W\,(\text{N})\\
\hline
(a) & 20.00 & 0.05 & 0.49 & 0.50 & 4.91\\
(b) & 0.50 & 2.00 & 19.62 & 20.00 & 196.20\\
(c) & 2.45 & 0.41 & 4.00 & 4.08 & 40.00\\
(d) & 0.10 & 10.00 & 98.10 & 100.00 & 981.00\\
(e) & 0.98 & 1.02 & 10.00 & 10.19 & 100.00
\end{array}
\]

\textbf{Final answer.}
\[
\boxed{\text{Completed table as above}}
\]

\textbf{Interpretation / common trap.} The main trap is unit consistency: litres must be converted to cubic metres before using $m=\rho V$.

\textbf{Source page.} Page 6.

\subsection*{Question 3}
\textbf{Rewritten question.} Determine the weight of a $12\,\text{kg}$ mass at a location where $g=9.77\,\text{m/s}^2$.

\textbf{Source wording caveat.} None.

\textbf{System type and boundary.} Closed system; weight is an external force acting on the mass.

\textbf{Assumptions.} Local gravity is constant over the object.

\textbf{Given / Find.}
Given: $m=12\,\text{kg}$, $g=9.77\,\text{m/s}^2$.
Find: $W$.

\textbf{Governing equation.}
\[
W=mg
\]

\textbf{Sequential substitution with units.}
\[
W=(12\,\text{kg})(9.77\,\text{m/s}^2)=117.24\,\text{N}
\]

\textbf{Final answer.}
\[
\boxed{W=117.24\,\text{N}}
\]

\textbf{Interpretation / common trap.} Weight is measured in newtons, not kilograms; do not leave the answer as a mass.

\textbf{Source page.} Page 7.

\subsection*{Question 4}
\textbf{Rewritten question.} Calculate the force required to accelerate an $8000\,\text{kg}$ rocket vertically upward at $28\,\text{m/s}^2$.

\textbf{Source wording caveat.} The solution deck labels the result as thrust force.

\textbf{System type and boundary.} Closed system for the rocket body; the force balance is external to the vehicle.

\textbf{Assumptions.} Gravity acts downward with $g=9.81\,\text{m/s}^2$; upward acceleration is uniform.

\textbf{Given / Find.}
Given: $m=8000\,\text{kg}$, $a=28\,\text{m/s}^2$, $g=9.81\,\text{m/s}^2$.
Find: $F_{\text{thrust}}$.

\textbf{Governing equation.}
\[
F_{\text{thrust}}=mg+ma=m(g+a)
\]

\textbf{Sequential substitution with units.}
\[
F_{\text{thrust}}=(8000\,\text{kg})(9.81\,\text{m/s}^2)+(8000\,\text{kg})(28\,\text{m/s}^2)
\]
\[
F_{\text{thrust}}=78{,}480\,\text{N}+224{,}000\,\text{N}=302{,}480\,\text{N}
\]

\textbf{Final answer.}
\[
\boxed{F_{\text{thrust}}=3.0248\times10^5\,\text{N}}
\]

\textbf{Interpretation / common trap.} The rocket must overcome both weight and provide the upward acceleration; forgetting the $mg$ term is the usual error.

\textbf{Source page.} Page 8.

\subsection*{Question 5}
\textbf{Rewritten question.} A $2200\,\text{kg}$ automobile travelling at $90\,\text{kph}$ collides with a stationary $1000\,\text{kg}$ automobile. After the collision, the large automobile slows to $60\,\text{kph}$ and the smaller vehicle moves at $78\,\text{kph}$. Determine the increase in internal energy of the two-car system.

\textbf{Source wording caveat.} None.

\textbf{System type and boundary.} Closed system containing both vehicles; collision is analysed through energy conservation.

\textbf{Assumptions.} Neglect external work and heat transfer during the short collision; all lost kinetic energy becomes internal energy of the system.

\textbf{Given / Find.}
Given: $m_A=2200\,\text{kg}$, $m_B=1000\,\text{kg}$, $V_{A1}=90\,\text{kph}$, $V_{A2}=60\,\text{kph}$, $V_{B1}=0$, $V_{B2}=78\,\text{kph}$.
Find: $\Delta U=U_2-U_1$.

\textbf{Governing equations.}
\[
\Delta U=E_{k1}-E_{k2},\qquad E_k=\frac12 mV^2
\]

\textbf{Sequential substitution with units.}
Convert speeds:
\[
90\,\text{kph}=25.0\,\text{m/s},\quad 60\,\text{kph}=16.67\,\text{m/s},\quad 78\,\text{kph}=21.67\,\text{m/s}
\]
Initial kinetic energy:
\[
E_{k1}=\tfrac12(2200\,\text{kg})(25.0\,\text{m/s})^2=687{,}550\,\text{J}
\]
Final kinetic energy:
\[
E_{k2}=\tfrac12(2200\,\text{kg})(16.67\,\text{m/s})^2+\tfrac12(1000\,\text{kg})(21.67\,\text{m/s})^2=540{,}472\,\text{J}
\]
Thus,
\[
\Delta U=687{,}550\,\text{J}-540{,}472\,\text{J}=147{,}078\,\text{J}
\]

\textbf{Final answer.}
\[
\boxed{\Delta U=1.47078\times10^5\,\text{J}}
\]

\textbf{Interpretation / common trap.} Use the kinetic energy of both cars after the collision; do not treat the stationary car as remaining stationary.

\textbf{Source page.} Page 9.

\subsection*{Question 6}
\textbf{Rewritten question.} Calculate the force due to pressure acting on a horizontal $1.2\,\text{m}$ diameter submarine hatch submerged $800\,\text{m}$ below the surface.

\textbf{Source wording caveat.} None.

\textbf{System type and boundary.} Hydrostatic pressure force on a flat circular hatch boundary.

\textbf{Assumptions.} Water density is $1000\,\text{kg/m}^3$; pressure variation over the hatch is neglected because only the gauge pressure at depth is used in the source solution.

\textbf{Given / Find.}
Given: $\rho=1000\,\text{kg/m}^3$, $h=800\,\text{m}$, $D=1.2\,\text{m}$.
Find: $F$.

\textbf{Governing equations.}
\[
p=\rho gh,\qquad A=\frac{\pi D^2}{4},\qquad F=pA
\]

\textbf{Sequential substitution with units.}
\[
p=(1000\,\text{kg/m}^3)(9.81\,\text{m/s}^2)(800\,\text{m})=7.848\times10^6\,\text{Pa}
\]
\[
A=\frac{\pi(1.2\,\text{m})^2}{4}=1.13097\,\text{m}^2
\]
\[
F=(7.848\times10^6\,\text{Pa})(1.13097\,\text{m}^2)=8.876\times10^6\,\text{N}
\]

\textbf{Final answer.}
\[
\boxed{F=8.876\times10^6\,\text{N}}
\]

\textbf{Interpretation / common trap.} Use the circular area with diameter $1.2\,\text{m}$, not radius; the result is a force, not pressure.

\textbf{Source page.} Page 10.

\subsection*{Question 7}
\questionfigure{week1-q07.png}{U-tube geometry with water, oil, and marked heights used in the pressure balance.}
\textbf{Rewritten question.} A U-tube with both arms open to the atmosphere contains water in one arm and water plus light oil in the other. Given $\rho_{oil}=790\,\text{kg/m}^3$ and $h_1=50\,\text{cm}$, determine the heights of the two fluid columns in the mixed-fluid arm.

\textbf{Source wording caveat.} The source figure is required to interpret the height labels. The solution deck indicates $h_3=h_2/4$ and uses that relation in the pressure balance.

\textbf{System type and boundary.} Static fluid columns; boundary is a level pressure comparison across the same horizontal datum.

\textbf{Assumptions.} Both free surfaces are exposed to atmosphere; the fluids are at rest; water density is $1000\,\text{kg/m}^3$.

\textbf{Given / Find.}
Given: $\rho_{water}=1000\,\text{kg/m}^3$, $\rho_{oil}=790\,\text{kg/m}^3$, $h_1=0.50\,\text{m}$, and $h_3=h_2/4$.
Find: $h_2$ and $h_3$.

\textbf{Governing equations.}
\[
p_{left}=p_{right},\qquad \rho_{water}gh_1-\rho_{water}gh_3=\rho_{oil}gh_2
\]
with $h_3=h_2/4$.

\textbf{Sequential substitution with units.}
\[
1000(0.50)-1000\left(\frac{h_2}{4}\right)=790h_2
\]
\[
500-250h_2=790h_2
\]
\[
500=1040h_2
\]
\[
h_2=\frac{500}{1040}=0.4808\,\text{m}
\]
\[
h_3=\frac{h_2}{4}=0.1202\,\text{m}
\]

\textbf{Final answer.}
\[
\boxed{h_2=0.48\,\text{m},\qquad h_3=0.12\,\text{m}}
\]

\textbf{Interpretation / common trap.} The key is matching points at the same elevation and keeping track of which fluid occupies each vertical segment.

\textbf{Source page.} Page 11.

\subsection*{Question 8}
\questionfigure{week1-q08.png}{Piston-cylinder with a compressed spring used to determine the gas pressure.}
\textbf{Rewritten question.} Calculate the pressure in a $150\,\text{mm}$ diameter cylinder when the piston has mass $40\,\text{kg}$ and the spring is compressed by $35\,\text{cm}$.

\textbf{Source wording caveat.} The solution deck does not restate the spring constant in words, but the working shows the spring force as $2000\times0.35$, so the recoverable value is $k=2000\,\text{N/m}$.

\textbf{System type and boundary.} Static force balance on the piston; boundary encloses the piston face.

\textbf{Assumptions.} Piston is in mechanical equilibrium; atmospheric pressure acts on the outside face; spring force acts upward/downward opposite the motion as shown in the source solution; the source result is reported as absolute pressure in the cylinder.

\textbf{Given / Find.}
Given: $m_p=40\,\text{kg}$, $k=2000\,\text{N/m}$, $x=0.35\,\text{m}$, $D=0.15\,\text{m}$.
Find: $p_{air}$.

\textbf{Governing equations.}
\[
pA=mg+kx,\qquad A=\frac{\pi D^2}{4}
\]

\textbf{Sequential substitution with units.}
\[
A=\frac{\pi(0.15\,\text{m})^2}{4}\approx0.0177\,\text{m}^2
\]
\[
mg=(40\,\text{kg})(9.81\,\text{m/s}^2)=392.4\,\text{N}
\]
\[
kx=(2000\,\text{N/m})(0.35\,\text{m})=700\,\text{N}
\]
\[
p=\frac{392.4\,\text{N}+700\,\text{N}}{0.0177\,\text{m}^2}=61{,}717.51\,\text{Pa}
\]

\textbf{Final answer.}
\[
\boxed{p=6.171751\times10^4\,\text{Pa}}
\]

\textbf{Interpretation / common trap.} Do not omit the spring force, and be careful that the piston area uses the diameter in metres.

\textbf{Source page.} Page 12.

\subsection*{Question 9}
\questionfigure{week1-q09.png}{Manometer geometry comparing water and oil columns with a mercury leg.}
\textbf{Rewritten question.} Determine the pressure difference between a water pipe and an oil pipe when the top of the mercury column is at the same elevation as the middle of the oil pipe.

\textbf{Source wording caveat.} The solution deck provides the geometry via a diagram, including the $6\,\text{m}$ mercury head and $3\,\text{m}$ water head.

\textbf{System type and boundary.} Static manometer balance.

\textbf{Assumptions.} Fluids are at rest; use $\rho_{Hg}=13600\,\text{kg/m}^3$ and $\rho_{water}=1000\,\text{kg/m}^3$ as in the source; the oil-side reference cancels in the provided working.

\textbf{Given / Find.}
Given: $h_{Hg}=6\,\text{m}$, $h_{water}=3\,\text{m}$.
Find: $\Delta p=p_{water}-p_{oil}$.

\textbf{Governing equations.}
\[
\Delta p=\rho_{Hg}gh_{Hg}-\rho_{water}gh_{water}
\]

\textbf{Sequential substitution with units.}
\[
\Delta p=(13600\,\text{kg/m}^3)(9.81\,\text{m/s}^2)(6\,\text{m})-(1000\,\text{kg/m}^3)(9.81\,\text{m/s}^2)(3\,\text{m})
\]
\[
\Delta p=800{,}496\,\text{Pa}-29{,}430\,\text{Pa}=771{,}066\,\text{Pa}
\]

\textbf{Final answer.}
\[
\boxed{\Delta p=7.71066\times10^5\,\text{Pa}}
\]

\textbf{Interpretation / common trap.} The pressure difference comes from vertical head differences, not from the horizontal pipe spacing.

\textbf{Source page.} Page 13.

\subsection*{Question 10}
\questionfigure{week1-q10.png}{Horizontal circular gate at the bottom of a water tank used for the opening-force calculation.}
\textbf{Rewritten question.} A horizontal circular gate of diameter $2.5\,\text{m}$ is located at the bottom of a water tank. Determine the force required to just open the gate.

\textbf{Source wording caveat.} The source page includes the gate geometry and a $6\,\text{m}$ water depth; the solution deck uses the moment balance about the hinge.

\textbf{System type and boundary.} Hydrostatic force and moment balance on a submerged gate.

\textbf{Assumptions.} Water density is $1000\,\text{kg/m}^3$; the resultant hydrostatic force acts at the gate center as shown in the source solution; the applied force is evaluated at the gate edge through a moment balance.

\textbf{Given / Find.}
Given: $D=2.5\,\text{m}$, $h=6\,\text{m}$, $\rho=1000\,\text{kg/m}^3$, $g=9.81\,\text{m/s}^2$.
Find: opening force $F$.

\textbf{Governing equations.}
\[
p=\rho gh,\qquad A=\frac{\pi D^2}{4},\qquad F_{water}=pA,
\qquad F\,D=F_{water}\,\frac{D}{2}
\]

\textbf{Sequential substitution with units.}
\[
p=(1000\,\text{kg/m}^3)(9.81\,\text{m/s}^2)(6\,\text{m})=58{,}860\,\text{Pa}
\]
\[
A=\frac{\pi(2.5\,\text{m})^2}{4}=4.9087\,\text{m}^2
\]
\[
F_{water}=(58{,}860\,\text{Pa})(4.9087\,\text{m}^2)=288{,}928.35\,\text{N}
\]
\[
F=\frac{F_{water}}{2}=144{,}464.17\,\text{N}
\]

\textbf{Final answer.}
\[
\boxed{F=1.44464\times10^5\,\text{N}}
\]

\textbf{Interpretation / common trap.} The opening force comes from a moment balance, not simply the full resultant hydrostatic force.

\textbf{Source page.} Page 14.

\weekstart{Week 2}{Energy, work, power and ideal-gas foundations}
\section*{Week 2}

\subsection*{Question 1}
\textbf{Rewritten question.} Determine the specific kinetic energy of a mass whose velocity is $45\,\mathrm{m/s}$, in $\mathrm{kJ/kg}$.

\textbf{Vagueness check.} No material ambiguity; the mass value is not needed because specific kinetic energy is per unit mass.

\textbf{System and boundary.} Closed body treated as a point mass; boundary is the mass itself.

\textbf{Assumptions.} Classical mechanics; no potential-energy change required.

\textbf{Given / Find.}
\[
V=45\,\mathrm{m/s},\qquad e_k=?
\]

\textbf{Governing equation.}
\[
 e_k=\frac{V^2}{2}
\]

\textbf{Substitution.}
\[
 e_k=\frac{(45\,\mathrm{m/s})^2}{2}=\frac{2025}{2}\,\mathrm{m^2/s^2}=1012.5\,\mathrm{J/kg}
\]
\[
 1012.5\,\mathrm{J/kg}=1.0125\,\mathrm{kJ/kg}
\]

\[\boxed{e_k=1.0125\,\mathrm{kJ/kg}}\]

\textbf{Reasonableness / common trap.} The common error is forgetting the factor of $1/2$ or reporting total kinetic energy instead of specific kinetic energy.

\textbf{Source page.} p.~1.

\subsection*{Question 2}
\textbf{Rewritten question.} A car is accelerated from rest to $85\,\mathrm{km/h}$ in $5\,\mathrm{s}$. Would the energy transferred to the car be different if it were accelerated to the same speed in $3.5\,\mathrm{s}$?

\textbf{Vagueness check.} The wording asks about energy transfer, not power. The mass is not given, so only a ratio or qualitative comparison is possible.

\textbf{System and boundary.} Car as a closed system; boundary encloses the vehicle. Energy crosses as work from the drivetrain.

\textbf{Assumptions.} Same final speed in both cases; start from rest; neglect losses, drag, rolling resistance, and elevation change.

\textbf{Given / Find.}
\[
V_0=0,\qquad V_1=85\,\mathrm{km/h},\qquad \Delta t_1=5\,\mathrm{s},\qquad \Delta t_2=3.5\,\mathrm{s}
\]
Find whether the transferred energy changes, and how the required average power compares.

\textbf{Governing equations.}
\[
W=\Delta KE=\frac12 m\left(V_1^2-V_0^2\right),
\qquad
\dot W_{\mathrm{avg}}=\frac{W}{\Delta t}
\]

\textbf{Substitution.}
Because $V_0=0$ and $V_1$ is the same in both cases, the ideal work input is identical:
\[
W_{5\,\mathrm{s}}=W_{3.5\,\mathrm{s}}
\]
For average power,
\[
\frac{\dot W_{3.5}}{\dot W_{5}}=\frac{W/3.5}{W/5}=\frac{5}{3.5}=1.4286
\]

\[\boxed{\text{Energy transferred is the same; }\dot W_{3.5}\approx1.43\,\dot W_{5}.}\]

\textbf{Reasonableness / common trap.} Shorter acceleration time does not change the kinetic-energy increase if the start and finish speeds are the same. The trap is confusing energy with power.

\textbf{Source page.} p.~2.

\subsection*{Question 3}
\textbf{Rewritten question.} A $5\,\mathrm{kg}$ system is heated from $50^\circ\mathrm{C}$ to $150^\circ\mathrm{C}$. Given $C_p=6000\,\mathrm{J/(kg\,K)}$, determine: (a) the heat added, and (b) the heat-transfer rate if the heating lasts $10\,\mathrm{s}$.

\textbf{Vagueness check.} Minor wording issue: the statement says $C_p$ but the process is treated as sensible heating with constant specific heat.

\textbf{System and boundary.} Closed system; boundary encloses the $5\,\mathrm{kg}$ system.

\textbf{Assumptions.} Constant specific heat; no phase change; negligible work.

\textbf{Given / Find.}
\[
m=5\,\mathrm{kg},\qquad c_p=6000\,\mathrm{J/(kg\,K)},\qquad T_1=50^\circ\mathrm{C},\qquad T_2=150^\circ\mathrm{C}
\]
Find $Q$ and $\dot Q$.

\textbf{Governing equations.}
\[
Q=mc_p\Delta T,\qquad \dot Q=\frac{Q}{\Delta t}
\]

\textbf{Substitution.}
\[
\Delta T=150-50=100\,\mathrm{K}
\]
\[
Q=(5\,\mathrm{kg})(6000\,\mathrm{J/(kg\,K)})(100\,\mathrm{K})=3.0\times10^6\,\mathrm{J}
\]
\[
Q=3000\,\mathrm{kJ}=3.00\,\mathrm{MJ}
\]
For $\Delta t=10\,\mathrm{s}$,
\[
\dot Q=\frac{3000\,\mathrm{kJ}}{10\,\mathrm{s}}=300\,\mathrm{kJ/s}=300\,\mathrm{kW}
\]

\[
\boxed{Q=3000\,\mathrm{kJ}},\qquad \boxed{\dot Q=300\,\mathrm{kW}}
\]

\textbf{Reasonableness / common trap.} Temperature rise is $100\,\mathrm{K}$, not $150\,\mathrm{K}$. Since $1\,\mathrm{kJ/s}=1\,\mathrm{kW}$, the rate is easy to misstate if units are not checked.

\textbf{Source page.} p.~3.

\subsection*{Question 4}
\textbf{Rewritten question.} A steel rod of diameter $0.8\,\mathrm{cm}$ and length $21\,\mathrm{m}$ is stretched $3\,\mathrm{cm}$. Young's modulus for this steel is $21\,\mathrm{kN/cm^2}$. How much work, in $\mathrm{kJ}$, is required to stretch this rod?

\textbf{Vagueness check.} The printed modulus appears inconsistent with steel. The literal printed value gives a very small answer; the likely intended value is $21{,}000\,\mathrm{kN/cm^2}$ (or $210\,\mathrm{GPa}$).

\textbf{System and boundary.} The rod is the system; the boundary is the rod surface. Work is done at the ends while the rod stretches elastically.

\textbf{Assumptions.} Linear elasticity; uniform axial stress; small strain; quasi-static loading.

\textbf{Given / Find.}
\[
d=0.8\,\mathrm{cm}=0.008\,\mathrm{m},\quad L=21\,\mathrm{m},\quad \delta=3\,\mathrm{cm}=0.03\,\mathrm{m},\quad E=21\,\mathrm{kN/cm^2}
\]
Find the required work.

\textbf{Governing equations.}
\[
A=\frac{\pi d^2}{4},\qquad \sigma=E\varepsilon,\qquad \varepsilon=\frac{\delta}{L},\qquad F=EA\frac{\delta}{L},\qquad W=\frac12F\delta=\frac{EA\delta^2}{2L}
\]

\textbf{Substitution.}
\[
A=\frac{\pi(0.008\,\mathrm{m})^2}{4}=5.0265\times10^{-5}\,\mathrm{m^2}
\]
Literal modulus conversion:
\[
1\,\mathrm{kN/cm^2}=10^7\,\mathrm{Pa}\quad\Rightarrow\quad E=21\times10^7=2.1\times10^8\,\mathrm{Pa}
\]
\[
F=(2.1\times10^8)(5.0265\times10^{-5})\frac{0.03}{21}=15.08\,\mathrm{N}
\]
\[
W=\frac12(15.08\,\mathrm{N})(0.03\,\mathrm{m})=0.226\,\mathrm{J}
\]
\[
W=2.262\times10^{-4}\,\mathrm{kJ}
\]
Likely intended steel modulus:
\[
E=21{,}000\,\mathrm{kN/cm^2}=2.1\times10^{11}\,\mathrm{Pa}
\]
This is $1000\times$ larger, so the work is also $1000\times$ larger:
\[
W\approx0.226\,\mathrm{kJ}
\]

\[
\boxed{W_{\text{literal}}=2.262\times10^{-4}\,\mathrm{kJ}},\qquad
\boxed{W_{\text{likely intended}}\approx0.226\,\mathrm{kJ}}
\]

\textbf{Reasonableness / common trap.} The trap is accepting the modulus literally without checking whether the steel value is plausible. The rod is very slender, so the force and work should be modest, but not absurdly tiny for steel unless the modulus is misprinted.

\textbf{Source page.} p.~4.

\subsection*{Question 5}
\textbf{Rewritten question.} Determine the torque applied to the shaft of a car that transmits $250\,\mathrm{hp}$ $(1\,\mathrm{hp}\approx745.7\,\mathrm{W})$ and rotates at a rate of $3200\,\mathrm{rpm}$.

\textbf{Vagueness check.} No ambiguity beyond the approximate horsepower conversion.

\textbf{System and boundary.} Rotating shaft as the control volume; power crosses the boundary as shaft work.

\textbf{Assumptions.} Steady rotation; no losses in the conversion from power to torque.

\textbf{Given / Find.}
\[
\dot W=250\,\mathrm{hp},\quad 1\,\mathrm{hp}=745.7\,\mathrm{W},\quad N=3200\,\mathrm{rpm},\quad \tau=?
\]

\textbf{Governing equation.}
\[
\dot W=\tau\omega
\]

\textbf{Substitution.}
\[
\dot W=(250)(745.7)=186425\,\mathrm{W}
\]
\[
\omega=3200\,\frac{\mathrm{rev}}{\mathrm{min}}\times\frac{2\pi\,\mathrm{rad}}{1\,\mathrm{rev}}\times\frac{1\,\mathrm{min}}{60\,\mathrm{s}}=335.103\,\mathrm{rad/s}
\]
\[
\tau=\frac{186425}{335.103}=556.32\,\mathrm{N\,m}
\]

\[\boxed{\tau\approx556.3\,\mathrm{N\,m}}\]

\textbf{Reasonableness / common trap.} A common trap is forgetting the $2\pi/60$ conversion from rpm to rad/s.

\textbf{Source page.} p.~5.

\subsection*{Question 6}
\textbf{Rewritten question.} A $100\,\mathrm{kW}$ engine and an $80\,\mathrm{MW}$ power plant are each used to lift a $500{,}000\,\mathrm{kg}$ mass $100\,\mathrm{m}$ vertically. Calculate the minimum time, in seconds, it would take for each power source to complete this task.

\textbf{Vagueness check.} No ambiguity if the problem is treated as an ideal lifting limit.

\textbf{System and boundary.} The lifted mass as the system; work crosses the boundary from the lifting device.

\textbf{Assumptions.} Ideal lifting, no losses, no acceleration losses, and constant $g$.

\textbf{Given / Find.}
\[
m=500000\,\mathrm{kg},\quad h=100\,\mathrm{m},\quad g=9.81\,\mathrm{m/s^2},\quad \dot W_1=100\,\mathrm{kW},\quad \dot W_2=80\,\mathrm{MW}
\]
Find $t_1$ and $t_2$.

\textbf{Governing equations.}
\[
W=mgh,\qquad t=\frac{W}{\dot W}
\]

\textbf{Substitution.}
\[
W=(500000)(9.81)(100)=4.905\times10^8\,\mathrm{J}
\]
For $100\,\mathrm{kW}$:
\[
t_1=\frac{4.905\times10^8\,\mathrm{J}}{100\times10^3\,\mathrm{W}}=4905\,\mathrm{s}
\]
For $80\,\mathrm{MW}$:
\[
t_2=\frac{4.905\times10^8\,\mathrm{J}}{80\times10^6\,\mathrm{W}}=6.13125\,\mathrm{s}
\]

\[
\boxed{t_1=4905\,\mathrm{s}},\qquad \boxed{t_2\approx6.13\,\mathrm{s}}
\]

\textbf{Reasonableness / common trap.} The times are ideal minimums; real systems take longer. Watch the megawatt/kilowatt conversion.

\textbf{Source page.} p.~6.

\subsection*{Question 7}
\questionfigure{week2-q07.png}{Vertical piston-cylinder water heating with heat loss and boundary work.}

\textbf{Rewritten question.} A vertical piston-cylinder device contains water and is being heated. During the process, $70\,\mathrm{kJ}$ of heat is transferred to the water, heat losses from the side walls amount to $9\,\mathrm{kJ}$, the piston rises as a result of evaporation, and $6\,\mathrm{kJ}$ of work is done by the vapor. Determine the change in the energy of the water for this process.

\textbf{Vagueness check.} The wording is clear once sign convention is chosen.

\textbf{System and boundary.} Closed system: the water in the cylinder. Boundary includes the piston and cylinder walls.

\textbf{Assumptions.} Closed system; heat into the system positive; work done by the system positive.

\textbf{Given / Find.}
\[
Q_{\text{in}}=70\,\mathrm{kJ},\quad Q_{\text{loss}}=9\,\mathrm{kJ},\quad W=6\,\mathrm{kJ}
\]
Find $\Delta E$.

\textbf{Governing equation.}
\[
\Delta E=Q-W
\]
with net heat
\[
Q=Q_{\text{in}}-Q_{\text{loss}}
\]

\textbf{Substitution.}
\[
Q=70-9=61\,\mathrm{kJ}
\]
\[
\Delta E=61-6=55\,\mathrm{kJ}
\]

\[\boxed{\Delta E=+55\,\mathrm{kJ}}\]

\textbf{Reasonableness / common trap.} The trap is double-counting the losses or flipping the work sign. Positive $\Delta E$ means the water's total energy increased.

\textbf{Source page.} p.~7.

\subsection*{Question 8}
\textbf{Rewritten question.} Consider a room that is initially at the outdoor temperature of $30^\circ\mathrm{C}$. The room contains a $50$-W lightbulb, a $120$-W TV set, a $300$-W refrigerator, and a $1200$-W iron. Assuming no heat transfer through the walls, determine the rate of increase of the energy content of the room when all of these electric devices are on.

\textbf{Vagueness check.} The initial temperature is irrelevant to the requested rate. The refrigerator still contributes positive electrical input to the room in the energy balance.

\textbf{System and boundary.} The room air and contents as a closed system; boundary is the room envelope.

\textbf{Assumptions.} No heat transfer through the walls; all electrical input eventually becomes energy increase in the room.

\textbf{Given / Find.}
\[
P_{\text{lamp}}=50\,\mathrm{W},\quad P_{\text{TV}}=120\,\mathrm{W},\quad P_{\text{fridge}}=300\,\mathrm{W},\quad P_{\text{iron}}=1200\,\mathrm{W}
\]
Find $\mathrm{d}E_{\text{room}}/\mathrm{d}t$.

\textbf{Governing equation.}
\[
\frac{\mathrm{d}E_{\text{room}}}{\mathrm{d}t}=\sum P_{\text{in}}
\]

\textbf{Substitution.}
\[
\frac{\mathrm{d}E_{\text{room}}}{\mathrm{d}t}=50+120+300+1200=1670\,\mathrm{W}
\]

\[
\boxed{\frac{\mathrm{d}E_{\text{room}}}{\mathrm{d}t}=1670\,\mathrm{W}=1.67\,\mathrm{kW}}\]

\textbf{Reasonableness / common trap.} Do not subtract the refrigerator as if it removed energy from the room overall; the electrical input still ends up as room energy.

\textbf{Source page.} p.~8.

\subsection*{Question 9}
\textbf{Rewritten question.} Consider a $2500\,\mathrm{kg}$ car cruising at constant speed of $80\,\mathrm{km/h}$. Now the car starts to pass another car by accelerating to $110\,\mathrm{km/h}$ in $5\,\mathrm{s}$. (a) Determine the additional power needed to achieve this acceleration. (b) What would your answer be if the total mass of the car were only $800\,\mathrm{kg}$?

\textbf{Vagueness check.} The phrase ``additional power'' is interpreted as the average rate associated with the kinetic-energy increase during the pass.

\textbf{System and boundary.} The car as a closed system.

\textbf{Assumptions.} Neglect losses; use only the change in kinetic energy; constant acceleration not required because only average power is asked.

\textbf{Given / Find.}
\[
m_1=2500\,\mathrm{kg},\quad m_2=800\,\mathrm{kg},\quad V_0=80\,\mathrm{km/h},\quad V_1=110\,\mathrm{km/h},\quad \Delta t=5\,\mathrm{s}
\]
Find additional average power for each mass.

\textbf{Governing equations.}
\[
W_a=\Delta KE=\frac12m\left(V_1^2-V_0^2\right),\qquad \dot W_a=\frac{W_a}{\Delta t}
\]

\textbf{Substitution.}
\[
V_0=\frac{80}{3.6}=22.2222\,\mathrm{m/s},\qquad V_1=\frac{110}{3.6}=30.5556\,\mathrm{m/s}
\]
\[
V_1^2-V_0^2=(30.5556)^2-(22.2222)^2=438.574\,\mathrm{m^2/s^2}
\]
For $m_1=2500\,\mathrm{kg}$,
\[
\dot W_{a,1}=\frac{1}{2}(2500)(438.574)\frac{1}{5}=1.09954\times10^5\,\mathrm{W}
\]
\[
\dot W_{a,1}\approx110\,\mathrm{kW}
\]
For $m_2=800\,\mathrm{kg}$,
\[
\dot W_{a,2}=\frac{1}{2}(800)(438.574)\frac{1}{5}=3.50859\times10^4\,\mathrm{W}
\]
\[
\dot W_{a,2}\approx35.1\,\mathrm{kW}
\]

\[
\boxed{\dot W_{a,1}\approx110\,\mathrm{kW}},\qquad \boxed{\dot W_{a,2}\approx35.1\,\mathrm{kW}}
\]

\textbf{Reasonableness / common trap.} The additional power scales linearly with mass. A common mistake is using speed in $\mathrm{km/h}$ directly in the kinetic-energy formula.

\textbf{Source page.} p.~9.

\subsection*{Question 10}
\questionfigure{week2-q10.png}{River flowing toward a lake, showing the 80 m elevation drop and river speed.}

\textbf{Rewritten question.} Consider a river flowing toward a lake at an average velocity of $4\,\mathrm{m/s}$ at a rate of $600\,\mathrm{m^3/s}$ at a location $80\,\mathrm{m}$ above the lake surface. Determine the total mechanical energy of the river water per unit mass and the power generation potential of the entire river at that location.

\textbf{Vagueness check.} Clear if the lake surface is used as the reference elevation and both free surfaces are at atmospheric pressure.

\textbf{System and boundary.} Flowing water stream through a control volume around the river section.

\textbf{Assumptions.} Incompressible stream; negligible pressure difference at the free surfaces; lake surface is the zero-elevation reference.

\textbf{Given / Find.}
\[
V=4\,\mathrm{m/s},\quad \dot V=600\,\mathrm{m^3/s},\quad z=80\,\mathrm{m},\quad p_1=p_2=0\;\text{(gauge)}
\]
Find the specific mechanical energy and the ideal power potential.

\textbf{Governing equations.}
\[
e_{mech}=\frac{p}{\rho}+\frac{V^2}{2}+gz
\]
For ideal generation potential,
\[
\dot m=\rho\dot V,\qquad \dot W_{gen}=\dot m\,e_{mech}
\]

\textbf{Substitution.}
Take the lake surface as $z=0$. At the river location,
\[
e_{mech}=\frac{0}{\rho}+\frac{(4\,\mathrm{m/s})^2}{2}+(9.81\,\mathrm{m/s^2})(80\,\mathrm{m})
\]
\[
e_{mech}=8+784.8=792.8\,\mathrm{J/kg}=0.7928\,\mathrm{kJ/kg}
\]
With $\rho\approx1000\,\mathrm{kg/m^3}$,
\[
\dot m=(1000)(600)=6.0\times10^5\,\mathrm{kg/s}
\]
\[
\dot W_{gen}=(6.0\times10^5)(792.8)=4.7568\times10^8\,\mathrm{W}
\]
\[
\dot W_{gen}\approx475.7\,\mathrm{MW}
\]

\[
\boxed{e_{mech}=0.7928\,\mathrm{kJ/kg}},\qquad \boxed{\dot W_{gen}\approx4.76\times10^8\,\mathrm{W}=475.7\,\mathrm{MW}}
\]

\textbf{Reasonableness / common trap.} The elevation term dominates the kinetic term. The usual trap is forgetting to convert the volumetric flow rate to mass flow rate.

\textbf{Source page.} p.~10.

\weekstart{Week 3}{Pure substances, interpolation and ideal gases}
% Week 3 – Thermo resource chapter
\section*{Week 3}
\addcontentsline{toc}{section}{Week 3}

\section*{Q1}
\textbf{Rewritten question.} For water at $100\,^{\circ}\mathrm{C}$, determine the saturation pressure and the saturated-liquid specific volume, internal energy, and enthalpy.

\textbf{Vague note.} Saturation-state lookup for water.

\textbf{System/boundary.} Closed pure-substance system at saturation; boundary is the water sample.

\textbf{Assumptions.}
\begin{itemize}
  \item Water is saturated liquid at $100\,^{\circ}\mathrm{C}$.
  \item Steam-table values are used directly.
\end{itemize}

\textbf{Given.}
\begin{itemize}
  \item $T = 100\,^{\circ}\mathrm{C}$
\end{itemize}

\textbf{Find.}
\begin{itemize}
  \item $P_{\mathrm{sat}}$, $v_f$, $u_f$, $h_f$
\end{itemize}

\textbf{Table choice.} Saturated water temperature table.

\textbf{Answer.}
\begin{itemize}
  \item $P_{\mathrm{sat}} = 101.42\ \mathrm{kPa}$
  \item $v_f = 0.001043\ \mathrm{m^3/kg}$
  \item $u_f = 419.06\ \mathrm{kJ/kg}$
  \item $h_f = 419.17\ \mathrm{kJ/kg}$
\end{itemize}

\[\boxed{P_{\mathrm{sat}} = 101.42\ \mathrm{kPa},\ v_f = 0.001043\ \mathrm{m^3/kg},\ u_f = 419.06\ \mathrm{kJ/kg},\ h_f = 419.17\ \mathrm{kJ/kg}}\]

\textbf{Trap.} Do not use atmospheric pressure; this is a saturation-table lookup.

\textbf{Source page.} week3-questions.pdf, Q1; official worked values from the local source page.

\section*{Q2}
\textbf{Rewritten question.} For R-134a at $500\ \mathrm{kPa}$, determine the saturation temperature and the saturated-vapour specific volume, internal energy, and enthalpy.

\textbf{Vague note.} Saturated R-134a state from pressure.

\textbf{System/boundary.} Closed pure-substance system at saturation.

\textbf{Assumptions.}
\begin{itemize}
  \item R-134a is saturated vapour at the stated pressure.
  \item Tables are read at the saturation pressure.
\end{itemize}

\textbf{Given.}
\begin{itemize}
  \item $P = 500\ \mathrm{kPa}$
\end{itemize}

\textbf{Find.}
\begin{itemize}
  \item $T_{\mathrm{sat}}$, $v_g$, $u_g$, $h_g$
\end{itemize}

\textbf{Table choice.} Saturated R-134a pressure table.

\textbf{Answer.}
\begin{itemize}
  \item $T_{\mathrm{sat}} = 15.71\,^{\circ}\mathrm{C}$
  \item $v_g = 0.041168\ \mathrm{m^3/kg}$
  \item $u_g = 238.77\ \mathrm{kJ/kg}$
  \item $h_g = 259.36\ \mathrm{kJ/kg}$
\end{itemize}

\[\boxed{T_{\mathrm{sat}} = 15.71\,^{\circ}\mathrm{C},\ v_g = 0.041168\ \mathrm{m^3/kg},\ u_g = 238.77\ \mathrm{kJ/kg},\ h_g = 259.36\ \mathrm{kJ/kg}}\]

\textbf{Trap.} Pressure is given, so use the saturation-pressure table, not the temperature table.

\textbf{Source page.} week3-questions.pdf, Q2; official worked values from the local source page.

\section*{Q3}
\textbf{Rewritten question.} Using linear interpolation, estimate the saturation pressure of water at $222.5\,^{\circ}\mathrm{C}$.

\textbf{Vague note.} Interpolate saturation pressure between nearby table entries.

\textbf{System/boundary.} Saturated water property lookup; no control-volume balance needed.

\textbf{Assumptions.}
\begin{itemize}
  \item Pressure varies linearly between the two nearest tabulated temperatures.
  \item Table values at $220\,^{\circ}\mathrm{C}$ and $225\,^{\circ}\mathrm{C}$ are accurate.
\end{itemize}

\textbf{Given.}
\begin{itemize}
  \item $T = 222.5\,^{\circ}\mathrm{C}$
  \item $(220\,^{\circ}\mathrm{C},\ 2319.6\ \mathrm{kPa})$
  \item $(225\,^{\circ}\mathrm{C},\ 2549.7\ \mathrm{kPa})$
\end{itemize}

\textbf{Find.}
\begin{itemize}
  \item $P_{\mathrm{sat}}$
\end{itemize}

\textbf{Table choice.} Saturated water temperature table plus interpolation.

\textbf{Answer.}
\begin{itemize}
  \item $P_{\mathrm{sat}} = 2434.65\ \mathrm{kPa}$
\end{itemize}

\[\boxed{P_{\mathrm{sat}} = 2434.65\ \mathrm{kPa}}\]

\textbf{Trap.} Interpolate on temperature, not pressure, and keep units in kPa.

\textbf{Source page.} week3-interpolation.pdf; official interpolated result provided by the user.

\section*{Q4}
\textbf{Rewritten question.} For superheated steam at $0.3\ \mathrm{MPa}$ and $325\,^{\circ}\mathrm{C}$, determine specific volume, internal energy, and enthalpy by interpolation.

\textbf{Vague note.} Superheated-steam interpolation at fixed pressure.

\textbf{System/boundary.} Closed steam sample in the superheated region.

\textbf{Assumptions.}
\begin{itemize}
  \item Linear interpolation in temperature is acceptable at $0.3\ \mathrm{MPa}$.
  \item The table values at $300\,^{\circ}\mathrm{C}$ and $400\,^{\circ}\mathrm{C}$ bracket the target state.
\end{itemize}

\textbf{Given.}
\begin{itemize}
  \item $P = 0.3\ \mathrm{MPa}$
  \item $T = 325\,^{\circ}\mathrm{C}$
  \item At $300\,^{\circ}\mathrm{C}$: $v=0.87535\ \mathrm{m^3/kg}$, $u=2807\ \mathrm{kJ/kg}$, $h=3069.6\ \mathrm{kJ/kg}$
  \item At $400\,^{\circ}\mathrm{C}$: $v=1.03155\ \mathrm{m^3/kg}$, $u=2966\ \mathrm{kJ/kg}$, $h=3275.5\ \mathrm{kJ/kg}$
\end{itemize}

\textbf{Find.}
\begin{itemize}
  \item $v$, $u$, $h$
\end{itemize}

\textbf{Table choice.} Superheated steam table at $0.3\ \mathrm{MPa}$ plus interpolation.

\textbf{Answer.}
\begin{itemize}
  \item $v = 0.9144\ \mathrm{m^3/kg}$
  \item $u = 2846.75\ \mathrm{kJ/kg}$
  \item $h = 3121.075\ \mathrm{kJ/kg}$
\end{itemize}

\[\boxed{v = 0.9144\ \mathrm{m^3/kg},\ u = 2846.75\ \mathrm{kJ/kg},\ h = 3121.075\ \mathrm{kJ/kg}}\]

\textbf{Trap.} Interpolate each property separately using the same temperature fraction.

\textbf{Source page.} week3-questions.pdf, Q4; official worked values from the local source page.

\section*{Q5}
\textbf{Rewritten question.} For saturated water at $90\,^{\circ}\mathrm{C}$ with specific volume $2\ \mathrm{m^3/kg}$, determine the quality.

\textbf{Vague note.} Quality from a two-phase specific-volume relation.

\textbf{System/boundary.} Closed saturated-water mixture.

\textbf{Assumptions.}
\begin{itemize}
  \item The state is a saturated mixture.
  \item $v = v_f + x(v_g-v_f)$ applies.
\end{itemize}

\textbf{Given.}
\begin{itemize}
  \item $T = 90\,^{\circ}\mathrm{C}$
  \item $v = 2\ \mathrm{m^3/kg}$
  \item $v_f = 0.001036\ \mathrm{m^3/kg}$
  \item $v_g = 2.3593\ \mathrm{m^3/kg}$
\end{itemize}

\textbf{Find.}
\begin{itemize}
  \item $x$
\end{itemize}

\textbf{Table choice.} Saturated water temperature table.

\textbf{Answer.}
\begin{itemize}
  \item $x = 0.847339$
\end{itemize}

\[\boxed{x = 0.847339}\]

\textbf{Trap.} The quality formula uses the saturated-mixture relation, not ideal-gas behavior.

\textbf{Source page.} week3-questions.pdf, Q5; official worked values from the local source page.

\section*{Q6}
\textbf{Rewritten question.} At $400\ \mathrm{kPa}$ and $200\,^{\circ}\mathrm{C}$, identify the phase of water and the correct property table. Then state the table to use if pressure stays at $400\ \mathrm{kPa}$ but temperature drops to $100\,^{\circ}\mathrm{C}$.

\textbf{Vague note.} Phase identification and table selection.

\textbf{System/boundary.} Water in a piston-cylinder device.

\textbf{Assumptions.}
\begin{itemize}
  \item Compare temperature to saturation temperature at the same pressure.
  \item Use standard steam tables for phase identification.
\end{itemize}

\textbf{Given.}
\begin{itemize}
  \item $P = 400\ \mathrm{kPa}$
  \item $T_1 = 200\,^{\circ}\mathrm{C}$
  \item $T_2 = 100\,^{\circ}\mathrm{C}$ at the same pressure
\end{itemize}

\textbf{Find.}
\begin{itemize}
  \item Phase at $200\,^{\circ}\mathrm{C}$ and the correct table
  \item Phase/table at $100\,^{\circ}\mathrm{C}$ and the correct table
\end{itemize}

\textbf{Table choice.} Compare against saturation temperature at $400\ \mathrm{kPa}$.

\textbf{Answer.}
\begin{itemize}
  \item At $400\ \mathrm{kPa}$, $T_{\mathrm{sat}} = 143.61\,^{\circ}\mathrm{C}$, so $200\,^{\circ}\mathrm{C}$ is superheated and the superheated steam table should be used.
  \item At $400\ \mathrm{kPa}$ and $100\,^{\circ}\mathrm{C}$, the water is compressed liquid and the compressed-liquid table should be used.
\end{itemize}

\[\boxed{\text{Superheated at }200\,^{\circ}\mathrm{C};\ \text{compressed liquid at }100\,^{\circ}\mathrm{C}}\]

\textbf{Trap.} The decision comes from $T$ versus $T_{\mathrm{sat}}$ at the same pressure, not from pressure alone.

\textbf{Source page.} week3-questions.pdf, Q6; official worked values from the local source page.

\section*{Q7}
\textbf{Rewritten question.} A rigid container holds $10\ \mathrm{kg}$ of water with total volume $1.115\ \mathrm{m^3}$ at $150\,^{\circ}\mathrm{C}$. It is heated until the pressure reaches $2000\ \mathrm{kPa}$. Determine the initial pressure, the initial quality, and the final temperature.

\textbf{Vague note.} Rigid-tank heating through the saturation dome into the superheated region.

\textbf{System/boundary.} Closed rigid container; no boundary work because volume is fixed.

\textbf{Assumptions.}
\begin{itemize}
  \item Mass is constant.
  \item Volume is constant.
  \item Initial state is a saturated mixture at $150\,^{\circ}\mathrm{C}$.
  \item Final state is superheated at $2\ \mathrm{MPa}$.
\end{itemize}

\textbf{Given.}
\begin{itemize}
  \item $m = 10\ \mathrm{kg}$
  \item $V = 1.115\ \mathrm{m^3}$
  \item $T_1 = 150\,^{\circ}\mathrm{C}$
  \item $P_2 = 2000\ \mathrm{kPa}$
\end{itemize}

\textbf{Find.}
\begin{itemize}
  \item $P_1$, $x_1$, $T_2$
\end{itemize}

\textbf{Table choice.} Saturated water temperature table at $150\,^{\circ}\mathrm{C}$, then superheated steam table at $2\ \mathrm{MPa}$.

\textbf{Answer.}
\begin{itemize}
  \item $v = V/m = 0.1115\ \mathrm{m^3/kg}$
  \item $P_1 = 476.16\ \mathrm{kPa}$
  \item $x_1 = 0.282095$
  \item $T_2 = 282.095\,^{\circ}\mathrm{C}$
\end{itemize}

\[\boxed{P_1 = 476.16\ \mathrm{kPa},\ x_1 = 0.282095,\ T_2 = 282.095\,^{\circ}\mathrm{C}}\]

\textbf{Trap.} Rigid container means use constant specific volume; do not apply a constant-pressure path to the first state.

\textbf{Source page.} week3-q1-q8.pdf, Q7; saturation and superheated-table extraction from the provided PDF/OCR, with the initial pressure read from the $150\,^{\circ}\mathrm{C}$ saturation entry.

\section*{Q8}
\textbf{Rewritten question.} A piston-cylinder device contains $0.6\ \mathrm{kg}$ of steam at $300\,^{\circ}\mathrm{C}$ and $0.5\ \mathrm{MPa}$. It is cooled at constant pressure until one-half of the mass condenses. Determine the final temperature and the volume change.

\textbf{Vague note.} Constant-pressure cooling into a saturated mixture.

\textbf{System/boundary.} Closed piston-cylinder system with moving boundary; pressure stays fixed.

\textbf{Assumptions.}
\begin{itemize}
  \item Pressure is constant throughout the process.
  \item Half the mass condenses, so final quality is $x_2 = 0.5$.
  \item Steam-table values are used for the initial superheated state and final saturated state.
\end{itemize}

\textbf{Given.}
\begin{itemize}
  \item $m = 0.6\ \mathrm{kg}$
  \item $P = 0.5\ \mathrm{MPa}$
  \item $T_1 = 300\,^{\circ}\mathrm{C}$
  \item Half the mass condenses $
    \Rightarrow x_2 = 0.5$
\end{itemize}

\textbf{Find.}
\begin{itemize}
  \item $T_2$
  \item $\Delta V = V_2 - V_1$
\end{itemize}

\textbf{Table choice.} Superheated steam table at $0.5\ \mathrm{MPa}$ for state 1, then saturated water table at $0.5\ \mathrm{MPa}$ for state 2.

\textbf{Answer.}
\begin{itemize}
  \item $T_2 = 151.83\,^{\circ}\mathrm{C}$
  \item $V_1 \approx 0.6\,(0.521\ \mathrm{m^3/kg}) \approx 0.313\ \mathrm{m^3}$
  \item $V_2 \approx 0.6\,[v_f + 0.5(v_g-v_f)] \approx 0.109\ \mathrm{m^3}$
  \item $\Delta V \approx -0.204\ \mathrm{m^3}$
\end{itemize}

\[\boxed{T_2 = 151.83\,^{\circ}\mathrm{C},\ \Delta V \approx -0.204\ \mathrm{m^3}}\]

\textbf{Trap.} The final temperature is the saturation temperature at $0.5\ \mathrm{MPa}$; the final quality is fixed by the condensation statement.

\textbf{Source page.} week3-q1-q8.pdf, Q8; final temperature from the saturation table at $0.5\ \mathrm{MPa}$. The volume-change figure is from the table-based state calculation, with OCR ambiguity in the source superheated entry.

\section*{Q9}
\textbf{Rewritten question.} A $3\ \mathrm{m^3}$ oxygen tank has a gauge pressure of $600\ \mathrm{kPa}$ at $28\,^{\circ}\mathrm{C}$. Determine the mass of oxygen.

\textbf{Vague note.} Ideal-gas mass from gauge pressure.

\textbf{System/boundary.} Closed gas tank.

\textbf{Assumptions.}
\begin{itemize}
  \item Oxygen behaves as an ideal gas.
  \item Absolute pressure is used.
\end{itemize}

\textbf{Given.}
\begin{itemize}
  \item $V = 3\ \mathrm{m^3}$
  \item $P_{\mathrm{gage}} = 600\ \mathrm{kPa}$
  \item $P_{\mathrm{atm}} = 97\ \mathrm{kPa}$
  \item $T = 28\,^{\circ}\mathrm{C} = 301.15\ \mathrm{K}$
\end{itemize}

\textbf{Find.}
\begin{itemize}
  \item $m_{\mathrm{O_2}}$
\end{itemize}

\textbf{Table choice.} Ideal-gas relation with oxygen gas constant.

\textbf{Answer.}
\begin{itemize}
  \item $P_{\mathrm{abs}} = 697\ \mathrm{kPa}$
  \item $m_{\mathrm{O_2}} = 26.7259\ \mathrm{kg}$
\end{itemize}

\[\boxed{m_{\mathrm{O_2}} = 26.7259\ \mathrm{kg}}\]

\textbf{Trap.} Use absolute pressure, not gauge pressure, in $Pv=mRT$.

\textbf{Source page.} week3-q9-q14.pdf, Q9; official worked values from the local source page.

\section*{Q10}
\textbf{Rewritten question.} An automobile tire of volume $0.5\ \mathrm{m^3}$ is inflated to $250\ \mathrm{kPa}$ gauge at $25\,^{\circ}\mathrm{C}$. Determine the mass of air.

\textbf{Vague note.} Ideal-gas mass from gauge pressure and standard atmosphere.

\textbf{System/boundary.} Closed air volume in the tire.

\textbf{Assumptions.}
\begin{itemize}
  \item Air behaves as an ideal gas.
  \item Standard atmospheric pressure applies.
\end{itemize}

\textbf{Given.}
\begin{itemize}
  \item $V = 0.5\ \mathrm{m^3}$
  \item $P_{\mathrm{gage}} = 250\ \mathrm{kPa}$
  \item $P_{\mathrm{atm}} = 101.325\ \mathrm{kPa}$
  \item $T = 25\,^{\circ}\mathrm{C} = 298.15\ \mathrm{K}$
\end{itemize}

\textbf{Find.}
\begin{itemize}
  \item $m_{\mathrm{air}}$
\end{itemize}

\textbf{Table choice.} Ideal-gas relation with air gas constant.

\textbf{Answer.}
\begin{itemize}
  \item $P_{\mathrm{abs}} = 351.325\ \mathrm{kPa}$
  \item $m_{\mathrm{air}} = 2.0529\ \mathrm{kg}$
\end{itemize}

\[\boxed{m_{\mathrm{air}} = 2.0529\ \mathrm{kg}}\]

\textbf{Trap.} Gauge pressure must be converted to absolute pressure before applying the ideal-gas equation.

\textbf{Source page.} week3-q9-q14.pdf, Q10; official worked values from the local source page.

\section*{Q11}
\textbf{Rewritten question.} A spherical balloon with diameter $10\ \mathrm{m}$ contains helium at $25\,^{\circ}\mathrm{C}$ and $250\ \mathrm{kPa}$ gauge. Determine the mole number and mass of helium.

\textbf{Vague note.} Ideal gas in a spherical container.

\textbf{System/boundary.} Closed helium balloon.

\textbf{Assumptions.}
\begin{itemize}
  \item The balloon is a perfect sphere.
  \item Helium behaves as an ideal gas.
  \item Absolute pressure is used.
\end{itemize}

\textbf{Given.}
\begin{itemize}
  \item $d = 10\ \mathrm{m}$
  \item $T = 25\,^{\circ}\mathrm{C} = 298.15\ \mathrm{K}$
  \item $P_{\mathrm{gage}} = 250\ \mathrm{kPa}$
  \item $P_{\mathrm{atm}} = 101.325\ \mathrm{kPa}$
\end{itemize}

\textbf{Find.}
\begin{itemize}
  \item $n_{\mathrm{He}}$, $m_{\mathrm{He}}$
\end{itemize}

\textbf{Table choice.} Ideal-gas relation with helium molecular data.

\textbf{Answer.}
\begin{itemize}
  \item $V = 523.5988\ \mathrm{m^3}$
  \item $P_{\mathrm{abs}} = 351.325\ \mathrm{kPa}$
  \item $m_{\mathrm{He}} = 297.069\ \mathrm{kg}$
  \item $n_{\mathrm{He}} \approx 74.21\ \mathrm{kmol}$
\end{itemize}

\[\boxed{m_{\mathrm{He}} = 297.069\ \mathrm{kg},\ n_{\mathrm{He}} \approx 74.21\ \mathrm{kmol}}\]

\textbf{Trap.} Use the balloon volume from sphere geometry before the gas law.

\textbf{Source page.} week3-q9-q14.pdf, Q11; official worked values from the local source page.

\section*{Q12}
\textbf{Rewritten question.} An ideal gas initially at $900\,^{\circ}\mathrm{C}$ expands from a partitioned rigid tank into vacuum, then is heated until the pressure returns to its initial value. Find the final temperature.

\textbf{Vague note.} Free expansion followed by constant-volume heating.

\textbf{System/boundary.} Closed ideal-gas system in a rigid tank.

\textbf{Assumptions.}
\begin{itemize}
  \item Ideal-gas behavior.
  \item The tank is rigid.
  \item Final pressure equals initial pressure.
\end{itemize}

\textbf{Given.}
\begin{itemize}
  \item $T_1 = 900\,^{\circ}\mathrm{C} = 1173.15\ \mathrm{K}$
  \item Expansion from one part to the full tank gives $V_2 = 4V_1$
  \item $P_3 = P_1$
\end{itemize}

\textbf{Find.}
\begin{itemize}
  \item $T_3$
\end{itemize}

\textbf{Table choice.} Ideal-gas relation between states.

\textbf{Answer.}
\begin{itemize}
  \item $T_3 = 4692.60\ \mathrm{K}$
\end{itemize}

\[\boxed{T_3 = 4692.60\ \mathrm{K}}\]

\textbf{Trap.} The free expansion changes volume, so the final heating must account for the $4\times$ total volume.

\textbf{Source page.} week3-q9-q14.pdf, Q12; official worked values from the local source page.

\section*{Q13}
\textbf{Rewritten question.} A $2\ \mathrm{m^3}$ tank of air at $10\,^{\circ}\mathrm{C}$ and $300\ \mathrm{kPa}$ is connected to a second tank that contains $3\ \mathrm{kg}$ of air at $35\,^{\circ}\mathrm{C}$ and $100\ \mathrm{kPa}$. After the valve is opened, the combined system reaches thermal equilibrium at $25\,^{\circ}\mathrm{C}$. Determine the second tank volume and the final absolute pressure.

\textbf{Vague note.} Two-tank ideal-gas mixing with final thermal equilibrium.

\textbf{System/boundary.} Combined closed system of both tanks.

\textbf{Assumptions.}
\begin{itemize}
  \item Air behaves as an ideal gas.
  \item Final temperature equals ambient temperature.
  \item Total mass is conserved.
\end{itemize}

\textbf{Given.}
\begin{itemize}
  \item Tank 1: $V_1 = 2\ \mathrm{m^3}$, $T_1 = 10\,^{\circ}\mathrm{C}$, $P_1 = 300\ \mathrm{kPa}$ gauge
  \item Tank 2: $m_2 = 3\ \mathrm{kg}$, $T_2 = 35\,^{\circ}\mathrm{C}$, $P_2 = 100\ \mathrm{kPa}$
  \item Final temperature: $T_f = 25\,^{\circ}\mathrm{C}$
\end{itemize}

\textbf{Find.}
\begin{itemize}
  \item $V_2$, $P_f$
\end{itemize}

\textbf{Table choice.} Ideal-gas equation of state for both tanks.

\textbf{Answer.}
\begin{itemize}
  \item $V_2 = 1.3179\ \mathrm{m^3}$
  \item $P_f = 332.1055\ \mathrm{kPa\ abs}$
\end{itemize}

\[\boxed{V_2 = 1.3179\ \mathrm{m^3},\ P_f = 332.1055\ \mathrm{kPa\ abs}}\]

\textbf{Trap.} The final pressure must be reported as absolute pressure.

\textbf{Source page.} week3-q9-q14.pdf, Q13; official worked values from the local source page.

\section*{Q14}
\textbf{Rewritten question.} Air in a piston-cylinder starts at $0.04\ \mathrm{m^3}$, $50\,^{\circ}\mathrm{C}$, and $500\ \mathrm{kPa}$ abs. Heat added is $60\ \mathrm{kJ}$ and the final temperature is $600\,^{\circ}\mathrm{C}$. With constant pressure and constant specific heats, determine the required paddle-wheel work added to the air.

\textbf{Vague note.} First-law balance with shaft work at constant pressure.

\textbf{System/boundary.} Closed air system in a piston-cylinder with paddle-wheel work.

\textbf{Assumptions.}
\begin{itemize}
  \item Air behaves as an ideal gas with constant specific heats.
  \item Pressure remains constant.
  \item The reported paddle-wheel work is work done on the air.
\end{itemize}

\textbf{Given.}
\begin{itemize}
  \item $V_1 = 0.04\ \mathrm{m^3}$
  \item $T_1 = 50\,^{\circ}\mathrm{C}$
  \item $P = 500\ \mathrm{kPa\ abs}$
  \item $Q_{\mathrm{in}} = 60\ \mathrm{kJ}$
  \item $T_2 = 600\,^{\circ}\mathrm{C}$
\end{itemize}

\textbf{Find.}
\begin{itemize}
  \item Paddle-wheel work input
\end{itemize}

\textbf{Table choice.} Ideal-gas energy relation with constant specific heats.

\textbf{Answer.}
\begin{itemize}
  \item $m = 0.2156\ \mathrm{kg}$
  \item $\Delta U = 119.199\ \mathrm{kJ}$
  \item Paddle-wheel work input $= 59.199\ \mathrm{kJ}$
\end{itemize}

\[\boxed{W_{\mathrm{paddle,in}} = 59.199\ \mathrm{kJ}}\] 

\textbf{Trap.} The arithmetic in the source has a small caveat; the energy balance still gives the stated result.

\textbf{Source page.} week3-q9-q14.pdf, Q14; official worked values from the local source page.
\weekstart{Week 4}{Boundary work and closed-system energy}
\section*{Week 4}

\subsection*{Q1}
\textbf{Rewritten question.} For the 3 kg nitrogen system shown on the P--v diagram, determine the total boundary work for the path $1\to 2\to 3$.

\textbf{Vague/source note.} The figure is diagram-dependent; the working page shows state 1 at $P_1=500\,\text{kPa}$, $v_1=0.5\,\text{m}^3/\text{kg}$, state 2 at $P_2=100\,\text{kPa}$, $v_2=1.0\,\text{m}^3/\text{kg}$, and state 3 at $P_3=400\,\text{kPa}$, $v_3=1.0\,\text{m}^3/\text{kg}$.

\textbf{System and boundary.} Closed system; nitrogen in a piston--cylinder style control mass. Boundary work is the area under the P--v path.

\textbf{Assumptions.} Quasi-equilibrium; sign convention $W_b>0$ for work done by the gas; process $2\to 3$ is constant volume.

\textbf{Given / Find.}
\begin{itemize}
  \item Given: $m=3\,\text{kg}$, $P_1=500\,\text{kPa}$, $v_1=0.5\,\text{m}^3/\text{kg}$, $P_2=100\,\text{kPa}$, $v_2=v_3=1.0\,\text{m}^3/\text{kg}$, $P_3=400\,\text{kPa}$.
  \item Find: $W_{1\to 3}$.
\end{itemize}

\textbf{Process / property model.} 
\[
W_{1\to 3}=m\int_{v_1}^{v_2} P\,dv + m\int_{v_2}^{v_3} P\,dv
\]
For the straight-line segment $1\to 2$, the source page evaluates the area as a triangle plus rectangle: $150\,\text{kJ/kg}$.
For $2\to 3$ at constant volume, $W_{2\to 3}=0$.

\textbf{Governing equations.}
\[
W_{1\to 2}=m\left[\frac{(v_2-v_1)(P_1-P_2)}{2}+(v_2-v_1)P_2\right],\qquad W_{2\to 3}=0
\]

\textbf{Source table values.} From the worked page: $W_{1\to 2}=150\,\text{kJ/kg}$.

\textbf{Sequential units.}
\[
W_{1\to 2,\,\text{total}}=(150\,\text{kJ/kg})(3\,\text{kg})=450\,\text{kJ}
\]
\[
W_{2\to 3,\,\text{total}}=0
\]
\[
W_{1\to 3}=450\,\text{kJ}
\]

\textbf{Boxed official answer.}
\[
\boxed{W_{1\to 3}=450\,\text{kJ}}
\]

\textbf{Trap.} Do not count any work for $2\to 3$; the vertical line means $dv=0$.

\textbf{Source pages.} Question sheet p. 1; worked page in \texttt{week4-q1-q4.pdf} p. 2.

\questionfigure{week4-q01.png}{Source $P$--$v$ path for Q1; boundary work is the signed area beneath the path.}

\subsection*{Q2}
\textbf{Rewritten question.} A $2\,\text{m}^3$ sample of saturated liquid water at $220^\circ\text{C}$ is expanded isothermally in a closed system until the final quality is 80\%. Determine the total work produced.

\textbf{Vague/source note.} The working page says the process remains inside the saturation dome, so the pressure stays at $P_{\text{sat}}(220^\circ\text{C})$.

\textbf{System and boundary.} Closed system; moving boundary work only.

\textbf{Assumptions.} Isothermal here means isobaric inside the saturation dome; quasi-equilibrium.

\textbf{Given / Find.}
\begin{itemize}
  \item Given: $V_1=2\,\text{m}^3$, $T=220^\circ\text{C}$, $x_2=0.8$.
  \item Find: $W_b$.
\end{itemize}

\textbf{Process / property model.} 
\[
W_b=P_{\text{sat}}(V_2-V_1),\qquad V_i=m v_i
\]
Using saturated-water values at $220^\circ\text{C}$, the source path treats the process as constant-pressure expansion.

\textbf{Governing equations.}
\[
m=\frac{V_1}{v_f},\qquad v_2=v_f+x_2(v_g-v_f),\qquad W_b=P_{\text{sat}}m(v_2-v_f)
\]

\textbf{Source table values.} Saturated water at $220^\circ\text{C}$: $P_{sat}=2319.6\,\text{kPa}$, $v_f=0.00119\,\text{m}^3/\text{kg}$ and $v_g=0.08609\,\text{m}^3/\text{kg}$.

\textbf{Sequential units.}
\[
 m=\frac{2}{0.00119}=1680.6723\,\text{kg}
\]
\[
 v_2=0.00119+0.8(0.08609-0.00119)=0.0691132\,\text{m}^3/\text{kg}
\]
\[
 w_b=2319.6(0.0691132-0.00119)=157.5547\,\text{kJ/kg}
\]
\[
 W_b=(1680.6723)(157.5547)=264797.74\,\text{kJ}
\]

\textbf{Boxed official answer.}
\[
\boxed{W_b\approx2.6480\times10^5\,\text{kJ}}
\]

\textbf{Trap.} Do not treat this as a fixed-temperature ideal-gas problem; the state lies in the two-phase region.

\textbf{Source pages.} Question sheet p. 1; worked page in \texttt{week4-q1-q4.pdf} p. 2.

\subsection*{Q3}
\textbf{Rewritten question.} Helium occupies $4\,\text{m}^3$ in a piston--cylinder device and is compressed at constant pressure $120\,\text{kPa}$ to $1\,\text{m}^3$. Find the initial temperature, final temperature, and compression work.

\textbf{Vague/source note.} The working page labels this as a constant-pressure compression process and uses the ideal-gas equation.

\textbf{System and boundary.} Closed system; moving boundary work.

\textbf{Assumptions.} Helium behaves as an ideal gas; pressure is constant.

\textbf{Given / Find.}
\begin{itemize}
  \item Given: $m=1\,\text{kg}$, $P=120\,\text{kPa}$, $V_1=4\,\text{m}^3$, $V_2=1\,\text{m}^3$.
  \item Find: $T_1$, $T_2$, $W_b$.
\end{itemize}

\textbf{Process / property model.}
\[
W_b=P(V_2-V_1),\qquad PV=mRT
\]
with $R_{\text{He}}=2.0769\,\mathrm{kJ/(kg\,K)}$.

\textbf{Governing equations.}
\[
T_1=\frac{P V_1}{mR},\qquad T_2=\frac{P V_2}{mR}
\]
\[
W_b=P(V_2-V_1)
\]

\textbf{Source table values.} Helium gas constant: $R=2.0769\,\mathrm{kJ/(kg\,K)}$.

\textbf{Sequential units.}
\[
T_1=\frac{120\times 4}{2.0769}=231.1\,\text{K}
\quad\text{and}\quad
T_2=\frac{120\times 1}{2.0769}=57.8\,\text{K}
\]
\[
W_b=120(1-4)=-360\,\text{kJ}
\]

\textbf{Boxed official answer.}
\[
\boxed{T_1\approx 231\,\text{K},\; T_2\approx 57.8\,\text{K},\; W_b=-360\,\text{kJ}}
\]

\textbf{Trap.} The negative sign means work is done on the gas during compression.

\textbf{Source pages.} Question sheet p. 1; worked page in \texttt{week4-q1-q4.pdf} p. 5.

\subsection*{Q4}
\textbf{Rewritten question.} A gas is compressed from $0.2\,\text{m}^3$ to $0.1\,\text{m}^3$ with $P=aV+b$, where $a=-1500\,\text{kPa/m}^3$ and $b=400\,\text{kPa}$. Find the work by geometry and by integration.

\textbf{Vague/source note.} Diagram-dependent path; no separate property table is needed.

\textbf{System and boundary.} Closed system; moving boundary work.

\textbf{Assumptions.} Quasi-equilibrium and the given linear pressure--volume relation.

\textbf{Given / Find.}
\begin{itemize}
  \item Given: $V_1=0.2\,\text{m}^3$, $V_2=0.1\,\text{m}^3$, $a=-1500\,\text{kPa/m}^3$, $b=400\,\text{kPa}$.
  \item Find: $W_b$.
\end{itemize}

\textbf{Process / property model.}
\[
P(V)=aV+b
\]

\textbf{Governing equations.}
\[
W_b=\int_{V_1}^{V_2}(aV+b)\,dV
=\frac{a}{2}(V_2^2-V_1^2)+b(V_2-V_1)
\]

\textbf{Source table values.} None.

\textbf{Sequential units.}
\[
W_b=\frac{-1500}{2}(0.1^2-0.2^2)+400(0.1-0.2)
=22.5-40=-17.5\,\text{kJ}
\]

\textbf{Boxed official answer.}
\[
\boxed{W_b=-17.5\,\text{kJ}}
\]

\textbf{Trap.} The path is a compression, so the boundary work must come out negative.

\textbf{Source pages.} Question sheet p. 1; worked page in \texttt{week4-q1-q4.pdf} p. 5.

\questionfigure{week4-q04.png}{Linear $P$--$V$ compression path for Q4; the negative signed area is work done on the gas.}

\subsection*{Q5}
\textbf{Rewritten question.} A $0.8\,\text{m}^3$ rigid tank contains R-134a at $180\,\text{kPa}$ and quality 45\%. Heat is added until the pressure reaches $600\,\text{kPa}$. Find the mass and heat transfer, and place the process on a $P$--$v$ diagram.

\textbf{Vague/source note.} The page is diagram-dependent; the worked page shows a constant-volume line from the saturated mixture state into the superheated region.

\textbf{System and boundary.} Closed rigid tank; no boundary work.

\textbf{Assumptions.} Fixed volume, negligible KE/PE changes, only heat transfer changes internal energy.

\textbf{Given / Find.}
\begin{itemize}
  \item Given: $V=0.8\,\text{m}^3$, $P_1=180\,\text{kPa}$, $x_1=0.45$, $P_2=600\,\text{kPa}$.
  \item Find: $m$, $Q$.
\end{itemize}

\textbf{Process / property model.}
\[
W_b=0,\qquad Q=\Delta U=m(u_2-u_1)
\]

\textbf{Governing equations.}
\[
v_1=v_f+x_1(v_g-v_f),\qquad m=\frac{V}{v_1},\qquad Q=m(u_2-u_1)
\]

\textbf{Source table values.} From Table A-12 at $180\,\text{kPa}$: $v_f=0.0007485\,\text{m}^3/\text{kg}$, $v_g=0.11049\,\text{m}^3/\text{kg}$, $u_f=34.81\,\text{kJ/kg}$, $u_{fg}=188.2\,\text{kJ/kg}$. Interpolated superheated value at $600\,\text{kPa}$ and constant $v$ gives $u_2\approx 327.2236\,\text{kJ/kg}$.

\textbf{Sequential units.}
\[
v_1=0.0007485+0.45(0.11049-0.0007485)=0.050132\,\text{m}^3/\text{kg}
\]
\[
m=\frac{0.8}{0.050132}=15.96\,\text{kg}
\]
\[
u_1=34.81+0.45(188.2)=119.5\,\text{kJ/kg}
\]
\[
Q=15.96(327.2236-119.5)=3315\,\text{kJ}
\]

\textbf{Boxed official answer.}
\[
\boxed{m=15.96\,\text{kg},\qquad Q\approx 3315\,\text{kJ}}
\]

\textbf{Trap.} Rigid tank means $W_b=0$ even though the pressure changes.

\textbf{Source pages.} Question sheet p. 1; worked pages in \texttt{week4-q5-q8.pdf} pp. 1--2.

\questionfigure{week4-q05.png}{Constant-volume heating path for the rigid R-134a tank in Q5.}

\subsection*{Q6}
\textbf{Rewritten question.} Steam in a spring-loaded piston--cylinder starts at $95\,\text{kPa}$ with quality 5\% and initial volume $1.2\,\text{m}^3$. It is heated until the volume is $5.5\,\text{m}^3$ and the pressure is $230\,\text{kPa}$. Find the heat transferred and work produced.

\textbf{Vague/source note.} The source page shows a linear spring relation on the $P$--$V$ diagram.

\textbf{System and boundary.} Closed system; moving boundary work against a spring.

\textbf{Assumptions.} Quasi-equilibrium, linear pressure--volume path, negligible KE/PE changes.

\textbf{Given / Find.}
\begin{itemize}
  \item Given: $P_1=95\,\text{kPa}$, $x_1=0.05$, $V_1=1.2\,\text{m}^3$, $P_2=230\,\text{kPa}$, $V_2=5.5\,\text{m}^3$.
  \item Find: $Q$, $W_b$.
\end{itemize}

\textbf{Process / property model.}
\[
W_b=\frac{P_1+P_2}{2}(V_2-V_1)
\]
\[
Q-W_b=\Delta U=m(u_2-u_1)
\]

\textbf{Governing equations.}
\[
m=\frac{V_1}{v_1},\qquad v_1=v_f+x_1(v_g-v_f)
\]
\[
Q=W_b+m(u_2-u_1)
\]

\textbf{Source table values.} From the steam tables at $95\,\text{kPa}$, the working page uses $v_f=0.001043\,\text{m}^3/\text{kg}$, $v_g=1.78696\,\text{m}^3/\text{kg}$, $u_f=384.36\,\text{kJ/kg}$, $u_g=2503.7\,\text{kJ/kg}$, giving $u_1=513.44\,\text{kJ/kg}$. At $230\,\text{kPa}$ the final state is in the two-phase region with $x_2\approx 0.5349$ and $u_2\approx 1598.05\,\text{kJ/kg}$.

\textbf{Sequential units.}
\[
v_1=0.001043+0.05(1.78696-0.001043)=0.0903377\,\text{m}^3/\text{kg}
\]
\[
m=\frac{1.2}{0.0903377}=13.28\,\text{kg}
\]
\[
W_b=\frac{95+230}{2}(5.5-1.2)=698.75\,\text{kJ}
\]
\[
Q=698.75+13.28(1598.05-513.44)=4698.75\,\text{kJ}
\]

\textbf{Boxed official answer.}
\[
\boxed{W_b\approx 698.75\,\text{kJ},\qquad Q\approx 4698.75\,\text{kJ}}
\]

\textbf{Trap.} The path is not constant pressure; use the straight-line spring work formula.

\textbf{Source pages.} Question sheet p. 1; worked pages in \texttt{week4-q5-q8.pdf} pp. 3--5.

\subsection*{Q7}
\textbf{Rewritten question.} Steel rods of diameter 10 cm, density $7833\,\text{kg/m}^3$, and $c_p=0.465\,\mathrm{kJ/(kg\,^{\circ}C)}$ move through a $1000^\circ\text{C}$ oven at $3\,\text{m/min}$, entering at $30^\circ\text{C}$ and leaving at $400^\circ\text{C}$. Find the heat-transfer rate to the rods.

\textbf{Vague/source note.} The working page treats this as a steady flow of material through the oven.

\textbf{System and boundary.} Steady-flow control volume around the oven section.

\textbf{Assumptions.} No work, no KE/PE changes, constant $c_p$.

\textbf{Given / Find.}
\begin{itemize}
  \item Given: $d=0.10\,\text{m}$, $\rho=7833\,\text{kg/m}^3$, $c_p=0.465\,\mathrm{kJ/(kg\,^{\circ}C)}$, $V=3\,\text{m/min}$, $T_1=30^\circ\text{C}$, $T_2=400^\circ\text{C}$.
  \item Find: $\dot Q$.
\end{itemize}

\textbf{Process / property model.}
\[
\dot Q=\dot m c_p(T_2-T_1),\qquad \dot m=\rho A V
\]

\textbf{Source table values.} Cross-sectional area $A=\pi(0.05)^2=0.00785\,\text{m}^2$.

\textbf{Sequential units.}
\[
\dot m=7833(0.00785)(3/60)=3.05\,\text{kg/s}
\]
\[
\dot Q=3.05(0.465)(400-30)=525\,\text{kW}
\]

\textbf{Boxed official answer.}
\[
\boxed{\dot Q\approx 525\,\text{kW}}
\]

\textbf{Trap.} Use the rods' mass flow rate, not the oven's volume flow rate alone.

\textbf{Source pages.} Question sheet p. 1; worked page in \texttt{week4-q5-q8.pdf} p. 6.

\subsection*{Q8}
\textbf{Rewritten question.} An egg of diameter 5.5 cm, density $1020\,\text{kg/m}^3$, and $c_p=3.32\,\mathrm{kJ/(kg\,^{\circ}C)}$ is heated from $4^\circ\text{C}$ to $90^\circ\text{C}$ in hot water at $97^\circ\text{C}$. Find the heat transferred.

\textbf{Vague/source note.} The egg is approximated as a uniform sphere.

\textbf{System and boundary.} Closed system; sensible heating only.

\textbf{Assumptions.} Uniform temperature in the egg; no phase change.

\textbf{Given / Find.}
\begin{itemize}
  \item Given: $d=5.5\,\text{cm}$, $\rho=1020\,\text{kg/m}^3$, $c_p=3.32\,\mathrm{kJ/(kg\,^{\circ}C)}$, $T_1=4^\circ\text{C}$, $T_2=90^\circ\text{C}$.
  \item Find: $Q$.
\end{itemize}

\textbf{Process / property model.}
\[
Q=mc_p\Delta T,\qquad m=\rho\frac{4}{3}\pi r^3
\]

\textbf{Source table values.} $r=0.0275\,\text{m}$.

\textbf{Sequential units.}
\[
V=\frac{4}{3}\pi(0.0275)^3=8.71\times 10^{-5}\,\text{m}^3
\]
\[
m=1020V=0.0888\,\text{kg}
\]
\[
Q=0.0888(3.32)(90-4)=25.35\,\text{kJ}
\]

\textbf{Boxed official answer.}
\[
\boxed{Q\approx 25.35\,\text{kJ}}
\]

\textbf{Trap.} Use the sphere volume, not a cylinder or shell approximation.

\textbf{Source pages.} Question sheet p. 1; worked page in \texttt{week4-q5-q8.pdf} p. 7.

\subsection*{Q9}
\textbf{Rewritten question.} Steam is held at constant pressure $400\,\text{kPa}$ in a piston--cylinder device. The initial state is $200^\circ\text{C}$ and $V_1=2\,\text{m}^3$. If $3500\,\text{kJ}$ of heat is added, determine the final temperature using enthalpy and the steam tables.

\textbf{Vague/source note.} The source pages confirm a constant-pressure enthalpy analysis; the final temperature comes from superheated-steam table interpolation.

\textbf{System and boundary.} Closed system with a frictionless piston.

\textbf{Assumptions.} Constant pressure, negligible KE/PE, steam tables apply.

\textbf{Given / Find.}
\begin{itemize}
  \item Given: $P=400\,\text{kPa}$, $T_1=200^\circ\text{C}$, $V_1=2\,\text{m}^3$, $Q=3500\,\text{kJ}$.
  \item Find: $T_2$.
\end{itemize}

\textbf{Process / property model.}
\[
Q=m(h_2-h_1)
\]
with $m=V_1/v_1$ and $h_2$ obtained from the $400\,\text{kPa}$ superheated table.

\textbf{Governing equations.}
\[
m=\frac{V_1}{v_1},\qquad h_2=h_1+\frac{Q}{m}
\]

\textbf{Source table values.} At $400\,\text{kPa}$ and $200^\circ\text{C}$, $v_1=0.53434\,\text{m}^3/\text{kg}$ and $h_1=2860.9\,\text{kJ/kg}$. At the same pressure, $h(600^\circ\text{C})=3703.3$ and $h(700^\circ\text{C})=3922.8\,\text{kJ/kg}$.

\textbf{Sequential units.}
\[
 m=\frac{2}{0.53434}=3.7429\,\text{kg},\qquad h_2=2860.9+\frac{3500}{3.7429}=3796.0038\,\text{kJ/kg}
\]
\[
 T_2=600+\frac{3796.0038-3703.3}{3922.8-3703.3}(700-600)=641.3303^\circ\text{C}
\]

\textbf{Boxed official answer.}
\[\boxed{T_2=641.33^\circ\text{C}}\]

\textbf{Trap.} Do not use saturation temperature at 400 kPa; the initial state is already superheated.

\textbf{Source pages.} Question sheet p. 2; worked pages in \texttt{week4-q9-q12.pdf} pp. 1, 13--14.

\subsection*{Q10}
\textbf{Rewritten question.} For superheated steam near $150\,\text{kPa}$, the specific heat is given by the supplied equation. Determine the enthalpy change from $400^\circ\text{C}$ to $800^\circ\text{C}$ for 4 kg of steam and the average value of $c_p$ over the interval.

\textbf{Vague/source note.} The supplied relation is $c_p(T)=2.07+(T-400)/1480$ in $\text{kJ/(kg}\cdot{}^\circ\text{C)}$, with $T$ in $^\circ\text{C}$.

\textbf{System and boundary.} Closed-system property change at constant pressure.

\textbf{Assumptions.} Superheated steam; the provided $c_p(T)$ relation is valid on the interval.

\textbf{Given / Find.}
\begin{itemize}
  \item Given: $m=4\,\text{kg}$, $T_1=400^\circ\text{C}$, $T_2=800^\circ\text{C}$, $P\approx 150\,\text{kPa}$.
  \item Find: $\Delta H$, $c_{p,\text{avg}}$.
\end{itemize}

\textbf{Process / property model.}
\[
\Delta h=\int_{T_1}^{T_2} c_p(T)\,dT,\qquad c_{p,\text{avg}}=\frac{1}{T_2-T_1}\int_{T_1}^{T_2} c_p(T)\,dT
\]

\textbf{Source table values.} Interpolation to $150\,\text{kPa}$ gives $h_1=3277.8\,\text{kJ/kg}$ at $400^\circ\text{C}$ and $h_2=4160.0\,\text{kJ/kg}$ at $800^\circ\text{C}$.

\textbf{Sequential units.}
\[
 \Delta H=4(4160.0-3277.8)=3528.8\,\text{kJ}
\]
Because the supplied $c_p(T)$ relation is linear,
\[
 c_p(400)=2.0700,\quad c_p(800)=2.3403,\quad c_{p,avg}=\frac{2.0700+2.3403}{2}=2.2052\,\text{kJ/(kg K)}.
\]

\textbf{Boxed official answer.}
\[\boxed{\Delta H=3528.8\,\text{kJ},\qquad c_{p,avg}=2.2052\,\text{kJ/(kg K)}}\]

\textbf{Trap.} The equation itself is part of the problem data; do not substitute a different correlation.

\textbf{Source pages.} Question sheet p. 2; worked pages in \texttt{week4-q9-q12.pdf} pp. 5--12.

\subsection*{Q11}
\textbf{Rewritten question.} Determine the enthalpy change for 2 kg of nitrogen heated from 250 K to 1300 K using (a) gas tables and (b) constant specific heat. Use $M=28\,\text{kg/kmol}$.

\textbf{Vague/source note.} The source table excerpt for ideal-gas properties is visible on the worked pages.

\textbf{System and boundary.} Ideal-gas property change.

\textbf{Assumptions.} Nitrogen behaves as an ideal gas; pressure does not affect $\Delta h$.

\textbf{Given / Find.}
\begin{itemize}
  \item Given: $m=2\,\text{kg}$, $T_1=250\,\text{K}$, $T_2=1300\,\text{K}$, $M=28\,\text{kg/kmol}$.
  \item Find: $\Delta H$ by gas tables and by constant $c_p$.
\end{itemize}

\textbf{Process / property model.}
\[
\Delta H=m(h_2-h_1)
\]

\textbf{Governing equations.}
\[
\Delta h = h_2-h_1\text{ from ideal-gas tables},\qquad \Delta h\approx c_p\Delta T\text{ for the constant-}c_p\text{ estimate}
\]

\textbf{Source table values.} The ideal-gas table difference, converted to mass basis, is $\Delta h=1175.1429\,\text{kJ/kg}$. The source's endpoint-average approximation uses $c_p(250)=1.039$ and $c_p(1300)=1.2215\,\text{kJ/(kg K)}$.

\textbf{Sequential units.}
\[
 \Delta H_{table}=2(1175.1429)=2350.2857\,\text{kJ}
\]
\[
 c_{p,avg}=\frac{1.039+1.2215}{2}=1.13025\,\text{kJ/(kg K)}
\]
\[
 \Delta H_{const\ c_p}=2(1.13025)(1300-250)=2373.53\,\text{kJ}
\]

\textbf{Boxed official answer.}
\[\boxed{\Delta H_{table}=2350.29\,\text{kJ},\qquad \Delta H_{const\ c_p}\approx2373.6\,\text{kJ}}\]

\textbf{Trap.} The tabulated enthalpy is molar; convert by dividing by molar mass.

\textbf{Source pages.} Question sheet p. 2; worked pages in \texttt{week4-q9-q12.pdf} pp. 13--14.

\subsection*{Q12}
\textbf{Rewritten question.} What is the enthalpy change, in kJ/kg, of oxygen as its temperature changes from 300 K to 400 K? Is there any difference if the temperature change were from 900 K to 1000 K?

\textbf{Vague/source note.} The answer is read from the ideal-gas property table for oxygen and then converted to mass basis.

\textbf{System and boundary.} Ideal-gas property change; pressure-independent enthalpy.

\textbf{Assumptions.} Oxygen behaves as an ideal gas.

\textbf{Given / Find.}
\begin{itemize}
  \item Given: $T_1=300\,\text{K}$, $T_2=400\,\text{K}$; compare with $900\,\text{K}$ to $1000\,\text{K}$.
  \item Find: $\Delta h$ per kg, and whether the interval matters.
\end{itemize}

\textbf{Process / property model.}
\[
\Delta h = \frac{h_2-h_1}{M}
\]
with $M_{\mathrm{O_2}}=31.999\,\text{kg/kmol}$.

\textbf{Source table values.} Table A-19 gives $\bar h_{300}=8736$, $\bar h_{400}=11711$, $\bar h_{900}=27928$ and $\bar h_{1000}=31389\,\text{kJ/kmol}$; $M_{O_2}=31.999\,\text{kg/kmol}$.

\textbf{Sequential units.}
\[
 \Delta h_{300\to400}=\frac{11711-8736}{31.999}=92.972\,\text{kJ/kg}
\]
\[
 \Delta h_{900\to1000}=\frac{31389-27928}{31.999}=108.160\,\text{kJ/kg}
\]

\textbf{Boxed official answer.}
\[\boxed{\Delta h_{300\to400}=92.972\,\text{kJ/kg},\quad \Delta h_{900\to1000}=108.160\,\text{kJ/kg}}\]

\textbf{Trap.} Do not use pressure in the ideal-gas enthalpy change; only temperature matters.

\textbf{Source pages.} Question sheet p. 2; worked pages in \texttt{week4-q9-q12.pdf} pp. 13--14.

\weekstart{Week 5}{Mass and energy analysis of control volumes}
\section*{Week 5}

\subsection*{Q1}
\textbf{Rewritten question.} An air compressor compresses $6\,\text{L}$ of air at $120\,\text{kPa}$ and $20^\circ\text{C}$ to $1000\,\text{kPa}$ and $400^\circ\text{C}$. Determine the flow work, in $\text{kJ/kg}$, required by the compressor.

\textbf{Vague note.} For an ideal gas, the specific flow work is $Pv=RT$; the compressor requirement is the change between the exit and inlet flow-work terms.

\textbf{Control-volume / system type.} Steady-flow compressor control volume; one inlet and one exit. Flow work crosses the boundary with the mass stream.

\textbf{Assumptions.} Air behaves as an ideal gas; kinetic and potential energy changes are irrelevant for flow work; steady operation.

\textbf{Given / Find.}
\begin{itemize}
  \item Given: $P_1=120\,\text{kPa}$, $T_1=293.15\,\text{K}$, $P_2=1000\,\text{kPa}$, $T_2=673.15\,\text{K}$, $R_{\text{air}}=0.2870\,\mathrm{kJ/(kg\,K)}$.
  \item Find: specific flow work required by the compressor.
\end{itemize}

\textbf{Mass balance.} For steady single-inlet/single-exit flow, $\dot m_{\text{in}}=\dot m_{\text{out}}$.

\textbf{Full SFEE then simplification.}
\[
\dot Q-\dot W_{\text{shaft}}+\dot m\left(h_1+\frac{V_1^2}{2}+gz_1\right)
=\dot m\left(h_2+\frac{V_2^2}{2}+gz_2\right)
\]
after canceling the negligible kinetic/potential terms and using ideal-gas flow work,
\[
w_{\text{flow, req}} = P_2v_2-P_1v_1 = R(T_2-T_1)
\]

\textbf{Table values.} Use the ideal-gas gas constant for air from Table A-1: $R=0.2870\,\mathrm{kJ/(kg\,K)}$.

\textbf{Sequential working with units.}
\[
\Delta T = 673.15-293.15 = 380.00\,\text{K}
\]
\[
w_{\text{flow, req}} = (0.2870\,\mathrm{kJ/(kg\,K)})(380.00\,\text{K}) = 109.06\,\text{kJ/kg}
\]

\textbf{Boxed result.}
\[
\boxed{w_{\text{flow, req}} = 109.06\,\text{kJ/kg}}
\]

\textbf{Trap.} Do not use the literal inlet volume $6\,\text{L}$ in the final specific flow-work answer; the requested quantity is per kilogram and depends on temperature via the ideal-gas relation.

\textbf{Source pages.} Question sheet p.~1; solution p.~1.

\subsection*{Q2}
\textbf{Rewritten question.} Air enters a $1\,\text{m}^2$ inlet of an aircraft engine at $100\,\text{kPa}$ and $20^\circ\text{C}$ with a velocity of $180\,\text{m/s}$. Determine the volume flow rate, in $\text{m}^3/\text{s}$, at the engine's inlet and the mass flow rate, in $\text{kg/s}$, at the engine's exit.

\textbf{Vague note.} The inlet volumetric flow rate comes from $\dot V = AV$; the exit mass flow rate is the same as the inlet mass flow rate for steady flow.

\textbf{Control-volume / system type.} Steady-flow engine control volume with one inlet and one exit.

\textbf{Assumptions.} Air is an ideal gas; steady operation; one-dimensional inlet; negligible accumulation inside the engine.

\textbf{Given / Find.}
\begin{itemize}
  \item Given: $A_1=1.0\,\text{m}^2$, $V_1=180\,\text{m/s}$, $P_1=100\,\text{kPa}$, $T_1=293.15\,\text{K}$, $R=0.2870\,\mathrm{kJ/(kg\,K)}$.
  \item Find: $\dot V_1$, $\dot m_2$.
\end{itemize}

\textbf{Mass balance.}
\[
\dot m_{\text{in}}=\dot m_{\text{out}}=\dot m
\]

\textbf{Full SFEE then simplification.} This problem is a mass-flow calculation, so the SFEE is not required beyond steady-flow continuity; the mass rate is enough.

\textbf{Table values.} For air as an ideal gas,
\[
\rho_1 = \frac{P_1}{RT_1}
\]
with $R=0.2870\,\mathrm{kJ/(kg\,K)} = 0.2870\,\mathrm{kPa\,m^3/(kg\,K)}$.

\textbf{Sequential working with units.}
\[
\dot V_1 = A_1V_1 = (1.0\,\text{m}^2)(180\,\text{m/s}) = 180\,\text{m}^3/\text{s}
\]
\[
\rho_1 = \frac{100\,\text{kPa}}{(0.2870\,\mathrm{kPa\,m^3/(kg\,K)})(293.15\,\text{K})} = 1.18858\,\text{kg/m}^3
\]
\[
\dot m_2=\dot m_1=\rho_1\dot V_1=(1.18858\,\text{kg/m}^3)(180\,\text{m}^3/\text{s})=213.94\,\text{kg/s}
\]

\textbf{Boxed result.}
\[
\boxed{\dot V_1 = 180\,\text{m}^3/\text{s}, \qquad \dot m_2 = 213.94\,\text{kg/s}}
\]

\textbf{Trap.} The exit mass flow rate is not found from the exit area here; the engine is a steady control volume, so the same mass rate passes through inlet and exit.

\textbf{Source pages.} Question sheet p.~1; solution p.~1.

\subsection*{Q3}
\textbf{Rewritten question.} Refrigerant-134a enters the compressor of a refrigeration system as saturated vapor at $0.14\,\text{MPa}$ and leaves as superheated vapor at $0.8\,\text{MPa}$ and $60^\circ\text{C}$ at a rate of $0.06\,\text{kg/s}$. Determine the rates of energy transfers by mass into and out of the compressor. Assume kinetic and potential energies are negligible.

\textbf{Vague note.} The compressor changes the refrigerant enthalpy; with negligible KE and PE, the mass-energy transfer rate is $\dot m h$ at each port.

\textbf{Control-volume / system type.} Steady-flow compressor control volume.

\textbf{Assumptions.} Steady operation; one inlet and one exit; negligible kinetic and potential energy changes; R-134a is a pure substance.

\textbf{Given / Find.}
\begin{itemize}
  \item Given: $\dot m=0.06\,\text{kg/s}$, inlet $P_1=0.14\,\text{MPa}$ saturated vapor, outlet $P_2=0.8\,\text{MPa}$, $T_2=60^\circ\text{C}$.
  \item Find: $\dot E_{\text{mass,in}}$, $\dot E_{\text{mass,out}}$.
\end{itemize}

\textbf{Mass balance.}
\[
\dot m_{\text{in}}=\dot m_{\text{out}}=0.06\,\text{kg/s}
\]

\textbf{Full SFEE then simplification.}
\[
\dot Q-\dot W_{\text{shaft}}+\dot m\left(h_1+\frac{V_1^2}{2}+gz_1\right)
=\dot m\left(h_2+\frac{V_2^2}{2}+gz_2\right)
\]
With negligible kinetic and potential energy changes, the rate of energy transfer by mass is simply
\[
\dot E_{\text{mass}}=\dot m h
\]

\textbf{Table values.} From the official solution pages:
\[
 h_1 = 239.19\,\text{kJ/kg}, \qquad h_2 = 296.82\,\text{kJ/kg}
\]

\textbf{Sequential working with units.}
\[
\dot E_{\text{mass,in}} = \dot m h_1 = (0.06\,\text{kg/s})(239.19\,\text{kJ/kg}) = 14.3514\,\text{kW}
\]
\[
\dot E_{\text{mass,out}} = \dot m h_2 = (0.06\,\text{kg/s})(296.82\,\text{kJ/kg}) = 17.8092\,\text{kW}
\]

\textbf{Boxed result.}
\[
\boxed{\dot E_{\text{mass,in}} = 14.3514\,\text{kW}, \qquad \dot E_{\text{mass,out}} = 17.8092\,\text{kW}}
\]

\textbf{Trap.} The compressor adds energy to the refrigerant; do not reverse the inlet and outlet enthalpies.

\textbf{Source pages.} Question sheet p.~1; solution p.~2--6.

\subsection*{Q4}
\textbf{Rewritten question.} Air enters a $16\,\text{cm}$-diameter pipe steadily at $200\,\text{kPa}$ and $20^\circ\text{C}$ with a velocity of $5\,\text{m/s}$. Air is heated as it flows, and it leaves the pipe at $180\,\text{kPa}$ and $40^\circ\text{C}$. Determine (a) the volume flow rate of air at the inlet, (b) the mass flow rate of air, and (c) the velocity and volume flow rate at the exit.

\textbf{Vague note.} Continuity gives the same mass flow rate through the pipe; ideal-gas density converts between $\dot V$ and $\dot m$.

\textbf{Control-volume / system type.} Steady-flow pipe control volume with one inlet and one exit.

\textbf{Assumptions.} Air is an ideal gas; steady flow; single inlet and exit; constant pipe diameter; negligible storage.

\textbf{Given / Find.}
\begin{itemize}
  \item Given: $D=0.16\,\text{m}$, $V_1=5\,\text{m/s}$, $P_1=200\,\text{kPa}$, $T_1=293.15\,\text{K}$, $P_2=180\,\text{kPa}$, $T_2=313.15\,\text{K}$.
  \item Find: (a) $\dot V_1$, (b) $\dot m$, (c) $V_2$ and $\dot V_2$.
\end{itemize}

\textbf{Mass balance.}
\[
\dot m_{\text{in}}=\dot m_{\text{out}}=\dot m
\]

\textbf{Full SFEE then simplification.} Not needed for the continuity calculation; only mass conservation and the ideal-gas relation are required.

\textbf{Table values.} Use $R_{\text{air}}=0.2870\,\mathrm{kJ/(kg\,K)}$.

\textbf{Sequential working with units.}
\[
A_1=\frac{\pi D^2}{4}=\frac{\pi(0.16\,\text{m})^2}{4}=0.0201062\,\text{m}^2
\]
\[
\dot V_1=A_1V_1=(0.0201062\,\text{m}^2)(5\,\text{m/s})=0.100531\,\text{m}^3/\text{s}
\]
\[
\rho_1=\frac{P_1}{RT_1}=\frac{200\,\text{kPa}}{(0.2870)(293\,\text{K})}=2.37838\,\text{kg/m}^3
\]
\[
\dot m=\rho_1\dot V_1=(2.37838\,\text{kg/m}^3)(0.100531\,\text{m}^3/\text{s})=0.23910\,\text{kg/s}
\]
At the exit, using $T_2=313\,\text{K}$,
\[
\rho_2=\frac{P_2}{RT_2}=\frac{180\,\text{kPa}}{(0.2870)(313\,\text{K})}=2.00376\,\text{kg/m}^3
\]
\[
\dot V_2=\frac{\dot m}{\rho_2}=0.119326\,\text{m}^3/\text{s}
\]
\[
V_2=\frac{\dot V_2}{A_1}=\frac{0.119326\,\text{m}^3/\text{s}}{0.0201062\,\text{m}^2}=5.93477\,\text{m/s}
\]

\textbf{Boxed result.}
\[
\boxed{\dot V_1 = 0.100531\,\text{m}^3/\text{s},\quad \dot m = 0.23910\,\text{kg/s},\quad \dot V_2 = 0.119326\,\text{m}^3/\text{s},\quad V_2 = 5.93477\,\text{m/s}}
\]

\textbf{Trap.} Do not forget the diameter must be converted to area before using $\dot V=AV$.

\textbf{Source pages.} Question sheet p.~2; solution p.~7.

\subsection*{Q5}
\textbf{Rewritten question.} A spherical hot-air balloon is initially filled with air at $120\,\text{kPa}$ and $20^\circ\text{C}$ with an initial diameter of $5\,\text{m}$. Air enters this balloon at $120\,\text{kPa}$ and $20^\circ\text{C}$ with a velocity of $3\,\text{m/s}$ through a $1\,\text{m}$-diameter opening. How many minutes will it take to inflate this balloon to a $17\,\text{m}$ diameter when the pressure and temperature of the air in the balloon remain the same as the air entering the balloon?

\textbf{Vague note.} Since the balloon pressure and temperature stay constant, the added mass is just the density times the change in balloon volume.

\textbf{Control-volume / system type.} Unsteady-flow control volume: the balloon is the control volume, and air mass accumulates inside.

\textbf{Assumptions.} Balloon is a perfect sphere; inlet flow is steady; balloon air remains at the inlet pressure and temperature; air is an ideal gas.

\textbf{Given / Find.}
\begin{itemize}
  \item Given: $D_1=5\,\text{m}$, $D_2=17\,\text{m}$, $P=120\,\text{kPa}$, $T=293.15\,\text{K}$, $V_{\text{in}}=3\,\text{m/s}$, inlet diameter $=1\,\text{m}$.
  \item Find: inflation time $t$.
\end{itemize}

\textbf{Mass balance.}
\[
\Delta m = m_2-m_1 = \dot m_{\text{in}}\,t
\]

\textbf{Full SFEE then simplification.} This is a filling problem, so only the unsteady mass balance is needed; the energy balance is not required to find the inflation time.

\textbf{Table values.} Use $R_{\text{air}}=0.2870\,\mathrm{kJ/(kg\,K)}$.

\textbf{Sequential working with units.}
\[
\rho = \frac{P}{RT} = \frac{120\,\text{kPa}}{(0.2870)(293.15\,\text{K})} = 1.42548\,\text{kg/m}^3
\]
\[
V_1=\frac{\pi}{6}(5\,\text{m})^3=65.4498\,\text{m}^3,\qquad
V_2=\frac{\pi}{6}(17\,\text{m})^3=2572.44\,\text{m}^3
\]
\[
A_{\text{in}}=\frac{\pi}{4}(1\,\text{m})^2=0.785398\,\text{m}^2
\]
\[
\dot V_{\text{in}}=A_{\text{in}}V_{\text{in}}=(0.785398)(3)=2.35619\,\text{m}^3/\text{s}
\]
\[
\dot m_{\text{in}}=\rho\dot V_{\text{in}}=(1.42548)(2.35619)=3.36063\,\text{kg/s}
\]
\[
\Delta m = \rho(V_2-V_1) = (1.42548)(2572.44-65.4498)=3575.71\,\text{kg}
\]
\[
t=\frac{\Delta m}{\dot m_{\text{in}}}=\frac{3575.71\,\text{kg}}{3.36063\,\text{kg/s}}=1064.6\,\text{s}=17.74\,\text{min}
\]

\textbf{Boxed result.}
\[
\boxed{t \approx 17.7\,\text{min}}
\]

\textbf{Trap.} The balloon volume change is large, but the answer is still based on mass added, not on pressure-volume work.

\textbf{Source pages.} Question sheet p.~2; solution p.~11.

\subsection*{Q6}
\textbf{Rewritten question.} A $2\,\text{m}^3$ rigid tank initially contains air whose density is $1.18\,\text{kg/m}^3$. The tank is connected to a high-pressure supply line through a valve. The valve is opened, and air is allowed to enter the tank until the density in the tank rises to $5.30\,\text{kg/m}^3$. Determine the mass of air that has entered the tank.

\textbf{Vague note.} Rigid tank means the volume does not change, so the mass added is just the final mass minus the initial mass.

\textbf{Control-volume / system type.} Unsteady rigid-tank control volume.

\textbf{Assumptions.} Tank volume is constant; density is uniform in the tank at each state; no leakage.

\textbf{Given / Find.}
\begin{itemize}
  \item Given: $V=2\,\text{m}^3$, $\rho_1=1.18\,\text{kg/m}^3$, $\rho_2=5.30\,\text{kg/m}^3$.
  \item Find: mass that entered, $m_{\text{in}}$.
\end{itemize}

\textbf{Mass balance.}
\[
 m_{\text{in}} = m_2-m_1
\]

\textbf{Full SFEE then simplification.} Energy balance is unnecessary for the requested mass gain; the rigid-tank volume and densities are sufficient.

\textbf{Table values.} None.

\textbf{Sequential working with units.}
\[
 m_1 = \rho_1V = (1.18\,\text{kg/m}^3)(2\,\text{m}^3)=2.36\,\text{kg}
\]
\[
 m_2 = \rho_2V = (5.30\,\text{kg/m}^3)(2\,\text{m}^3)=10.60\,\text{kg}
\]
\[
 m_{\text{in}} = 10.60-2.36 = 8.24\,\text{kg}
\]

\textbf{Boxed result.}
\[
\boxed{m_{\text{in}} = 8.24\,\text{kg}}
\]

\textbf{Trap.} Do not use the supply-line conditions; only the tank densities and fixed volume are needed.

\textbf{Source pages.} Question sheet p.~2; solution p.~13.

\subsection*{Q7}
\textbf{Rewritten question.} Air at $600\,\text{kPa}$ and $500\,\text{K}$ enters an adiabatic nozzle that has an inlet-to-exit area ratio of $2{:}1$ with a velocity of $120\,\text{m/s}$ and leaves with a velocity of $380\,\text{m/s}$. Determine (a) the exit temperature and (b) the exit pressure of the air.

\textbf{Vague note.} In an adiabatic nozzle, the inlet enthalpy drop pays for the rise in kinetic energy.

\textbf{Control-volume / system type.} Steady-flow nozzle control volume.

\textbf{Assumptions.} Adiabatic; no shaft work; negligible potential-energy change; air is an ideal gas with temperature-dependent enthalpy.

\textbf{Given / Find.}
\begin{itemize}
  \item Given: $P_1=600\,\text{kPa}$, $T_1=500\,\text{K}$, $V_1=120\,\text{m/s}$, $V_2=380\,\text{m/s}$, $A_1/A_2=2$.
  \item Find: $T_2$, $P_2$.
\end{itemize}

\textbf{Mass balance.}
\[
\dot m_{\text{in}}=\dot m_{\text{out}}
\]

\textbf{Full SFEE then simplification.}
\[
\dot Q-\dot W_{\text{shaft}}+\dot m\left(h_1+\frac{V_1^2}{2}+gz_1\right)
=\dot m\left(h_2+\frac{V_2^2}{2}+gz_2\right)
\]
With $\dot Q=0$, $\dot W_{\text{shaft}}=0$, and negligible $gz$ terms,
\[
h_1+\frac{V_1^2}{2}=h_2+\frac{V_2^2}{2}
\]

\textbf{Table values.} Use Table A-17 for air enthalpy at $T_1=500\,\text{K}$. For the pressure relation, use the ideal-gas law and continuity.

\textbf{Sequential working with units.}
\[
h_2=h_1+\frac{V_1^2-V_2^2}{2}
\]
With $c_p\approx 1.005\,\mathrm{kJ/(kg\,K)}$ and the standard air-table value at $500\,\text{K}$, this gives
\[
T_2 \approx 435.3\,\text{K}
\]
The density ratio from continuity and the area ratio gives
\[
\frac{\rho_2}{\rho_1}=\frac{A_1V_1}{A_2V_2}=\frac{2\times 120}{380}=0.63158
\]
Using $\rho\propto P/T$ for an ideal gas,
\[
\frac{P_2}{P_1}=\frac{\rho_2}{\rho_1}\frac{T_2}{T_1}
=0.63158\left(\frac{435.3}{500}\right)
\]
\[
P_2 \approx 3.30\times 10^2\,\text{kPa}
\]

\textbf{Boxed result.}
\[
\boxed{T_2 \approx 435\,\text{K}, \qquad P_2 \approx 330\,\text{kPa}}
\]

\textbf{Trap.} Do not use a constant-$c_p$ shortcut without checking the outlet pressure with continuity and the area ratio.

\textbf{Source pages.} Question sheet p.~3; solution p.~14.

\subsection*{Q8}
\textbf{Rewritten question.} An adiabatic air compressor compresses $10\,\text{L/s}$ of air at $120\,\text{kPa}$ and $20^\circ\text{C}$ to $1000\,\text{kPa}$ and $300^\circ\text{C}$. Determine (a) the work required by the compressor, in $\text{kJ/kg}$, and (b) the power required to drive the air compressor, in $\text{kW}$.

\textbf{Vague note.} For an adiabatic compressor, the shaft work input equals the enthalpy rise for an ideal-gas model.

\textbf{Control-volume / system type.} Steady-flow compressor control volume.

\textbf{Assumptions.} Adiabatic; negligible KE and PE changes; air is an ideal gas; single inlet and exit.

\textbf{Given / Find.}
\begin{itemize}
  \item Given: $\dot V_1=10\,\text{L/s}=0.010\,\text{m}^3/\text{s}$, $P_1=120\,\text{kPa}$, $T_1=293.15\,\text{K}$, $P_2=1000\,\text{kPa}$, $T_2=573.15\,\text{K}$.
  \item Find: $w_{\text{in}}$, $\dot W_{\text{in}}$.
\end{itemize}

\textbf{Mass balance.}
\[
\dot m_{\text{in}}=\dot m_{\text{out}}=\dot m
\]

\textbf{Full SFEE then simplification.}
\[
\dot Q-\dot W_{\text{shaft}}+\dot m\left(h_1+\frac{V_1^2}{2}+gz_1\right)
=\dot m\left(h_2+\frac{V_2^2}{2}+gz_2\right)
\]
With $\dot Q=0$ and negligible KE/PE,
\[
 w_{\text{in}} = h_2-h_1
\]

\textbf{Table values.} From the air table, interpolate at $T_1=293.15\,\text{K}$ and read the value at $T_2=573.15\,\text{K}$:
\[
 h_1 \approx 293.32\,\text{kJ/kg}, \qquad h_2 \approx 578.88\,\text{kJ/kg}
\]

\textbf{Sequential working with units.}
\[
 w_{\text{in}} = h_2-h_1 = 578.88-293.32 = 285.56\,\text{kJ/kg}
\]
\[
\rho_1=\frac{P_1}{RT_1}=\frac{120}{(0.2870)(293.15)}=1.4263\,\text{kg/m}^3
\]
\[
\dot m = \rho_1\dot V_1=(1.4263)(0.010)=0.01426\,\text{kg/s}
\]
\[
\dot W_{\text{in}} = \dot m\,w_{\text{in}} = (0.01426\,\text{kg/s})(285.56\,\text{kJ/kg}) = 4.08\,\text{kW}
\]

\textbf{Boxed result.}
\[
\boxed{w_{\text{in}} = 285.56\,\text{kJ/kg}, \qquad \dot W_{\text{in}} = 4.08\,\text{kW}}
\]

\textbf{Trap.} The compressor work is an input, so the sign convention should not be reported as negative when the question asks for required work.

\textbf{Source pages.} Question sheet p.~3; solution p.~20.

\subsection*{Q9}
\textbf{Rewritten question.} Refrigerant-134a is throttled from the saturated liquid state at $700\,\text{kPa}$ to a pressure of $160\,\text{kPa}$. Determine the temperature drop during this process and the final specific volume of the refrigerant.

\textbf{Vague note.} Throttling is an isenthalpic process, so the exit state is found by holding $h$ constant and then reading the mixture properties.

\textbf{Control-volume / system type.} Steady-flow throttling valve.

\textbf{Assumptions.} Adiabatic; no shaft work; negligible kinetic and potential energy changes; R-134a is a pure substance.

\textbf{Given / Find.}
\begin{itemize}
  \item Given: inlet saturated liquid at $700\,\text{kPa}$, exit pressure $160\,\text{kPa}$.
  \item Find: $\Delta T$ and $v_2$.
\end{itemize}

\textbf{Mass balance.}
\[
\dot m_{\text{in}}=\dot m_{\text{out}}
\]

\textbf{Full SFEE then simplification.}
\[
\dot Q-\dot W_{\text{shaft}}+\dot m\left(h_1+\frac{V_1^2}{2}+gz_1\right)
=\dot m\left(h_2+\frac{V_2^2}{2}+gz_2\right)
\]
For a throttling valve,
\[
 h_1=h_2
\]

\textbf{Table values.} From the official solution pages:
\[
T_1=26.69^\circ\text{C}, \qquad h_1=88.82\,\text{kJ/kg}
\]
At $160\,\text{kPa}$, saturation gives
\[
T_2=15.6^\circ\text{C}
\]
and with the quality computed from the saturated properties,
\[
x=0.2712,
\quad v_f=0.0007435\,\text{m}^3/\text{kg},
\quad v_g=0.12355\,\text{m}^3/\text{kg}
\]

\textbf{Sequential working with units.}
\[
\Delta T = T_1-T_2 = 26.69-15.6 = 11.09^\circ\text{C}
\]
\[
v_2=v_f+x(v_g-v_f)
=(0.0007435)+(0.2712)(0.12355-0.0007435)
=0.0340\,\text{m}^3/\text{kg}
\]

\textbf{Boxed result.}
\[
\boxed{\Delta T = 11.09^\circ\text{C}, \qquad v_2 = 0.0340\,\text{m}^3/\text{kg}}
\]

\textbf{Trap.} The exit is a vapour--liquid mixture, not superheated vapour; using saturated-mixture relations is the key step.

\textbf{Source pages.} Question sheet p.~4; solution p.~22.

\subsection*{Q10}
\textbf{Rewritten question.} Refrigerant-134a at $1\,\text{MPa}$ and $90^\circ\text{C}$ is to be cooled to $1\,\text{MPa}$ and $30^\circ\text{C}$ in a condenser by air. The air enters at $100\,\text{kPa}$ and $27^\circ\text{C}$ with a volume flow rate of $600\,\text{m}^3/\text{min}$ and leaves at $95\,\text{kPa}$ and $60^\circ\text{C}$. Determine the mass flow rate of the refrigerant.

\textbf{Vague note.} The heat lost by the refrigerant equals the heat gained by the air.

\textbf{Control-volume / system type.} Steady-flow heat exchanger with two control volumes: one for R-134a and one for air.

\textbf{Assumptions.} No work transfer; negligible KE and PE changes; R-134a is a pure substance; air is an ideal gas.

\textbf{Given / Find.}
\begin{itemize}
  \item Given: R-134a at $P_3=1\,\text{MPa}$, $T_3=90^\circ\text{C}$, $P_4=1\,\text{MPa}$, $T_4=30^\circ\text{C}$; air at $P_1=100\,\text{kPa}$, $T_1=27^\circ\text{C}=300\,\text{K}$, $\dot V_1=600\,\text{m}^3/\text{min}=10\,\text{m}^3/\text{s}$, $P_2=95\,\text{kPa}$, $T_2=60^\circ\text{C}=333\,\text{K}$.
  \item Find: $\dot m_{\text{R134a}}$.
\end{itemize}

\textbf{Mass balance.}
\[
\dot m_{\text{air,in}}=\dot m_{\text{air,out}},\qquad
\dot m_{\text{R,in}}=\dot m_{\text{R,out}}=\dot m_{\text{R134a}}
\]

\textbf{Full SFEE then simplification.}
For each stream in the heat exchanger,
\[
\dot Q-\dot W_{\text{shaft}}+\dot m\left(h_1+\frac{V_1^2}{2}+gz_1\right)
=\dot m\left(h_2+\frac{V_2^2}{2}+gz_2\right)
\]
With $\dot W_{\text{shaft}}=0$ and negligible KE/PE changes,
\[
\dot Q = \dot m(h_2-h_1)
\]
For the coupled heat exchanger,
\[
\dot Q_{\text{air}} = -\dot Q_{\text{R134a}}
\]

\textbf{Table values.} For air, use $R=0.2870\,\mathrm{kJ/(kg\,K)}$ and the air enthalpy values from Table A-17 at $300\,\text{K}$ and $333\,\text{K}$. For R-134a, use the official table values from the solution pages:
\[
 h_3 \approx 416.03\,\text{kJ/kg}, \qquad h_4 \approx 218.19\,\text{kJ/kg}
\]
so that
\[
 h_3-h_4 \approx 197.84\,\text{kJ/kg}
\]

\textbf{Sequential working with units.}
\[
\rho_{\text{air},1}=\frac{P_1}{RT_1}=\frac{100\,\text{kPa}}{(0.2870)(300\,\text{K})}=1.16144\,\text{kg/m}^3
\]
\[
\dot m_{\text{air}}=\rho_{\text{air},1}\dot V_1=(1.16144\,\text{kg/m}^3)(10\,\text{m}^3/\text{s})=11.6144\,\text{kg/s}
\]
Using the air-table enthalpy rise,
\[
\dot Q_{\text{air}}=\dot m_{\text{air}}(h_2-h_1)\approx 115.6\,\text{kW}
\]
Therefore, for the refrigerant,
\[
\dot Q_{\text{air}} = \dot m_{\text{R134a}}(h_3-h_4)
\]
\[
\dot m_{\text{R134a}} = \frac{115.6\,\text{kW}}{197.84\,\text{kJ/kg}} = 0.584\,\text{kg/s}
\]

\textbf{Boxed result.}
\[
\boxed{\dot m_{\text{R134a}} \approx 0.584\,\text{kg/s}}
\]

\textbf{Trap.} The air-side heat gain must be computed first; do not try to solve the refrigerant stream alone because the exchanger couples both flows.

\textbf{Source pages.} Question sheet p.~4; solution p.~23.

\weekstart{Week 6}{Second law, engines and refrigeration}
\section*{Week 6}

\subsection*{Q1}
\textbf{Rewritten question.} A steam power plant delivers an actual power output of $150\,\text{MW}$. It consumes coal at a rate of $60\,\text{tons/h}$, and the coal heating value is $30{,}000\,\text{kJ/kg}$. Determine the overall efficiency of the plant.

\textbf{Vague note.} "Overall efficiency" means actual useful power out divided by the ideal thermal power available from the fuel.

\textbf{Device/system type and reservoir boundary.} Steady-flow power plant; treat the plant as the control volume. Fuel energy enters from the hot reservoir side, and useful power leaves at the shaft/electrical output boundary.

\textbf{Assumptions.} Coal mass flow is uniform; heating value is fully available as thermal input; moisture and other losses are neglected for the ideal thermal input calculation; steady operation.

\textbf{Given / Find.}
\begin{itemize}
  \item Given: $\dot W_{\text{out}}=150\,\text{MW}$, coal rate $=60\,\text{tons/h}$, heating value $=30{,}000\,\text{kJ/kg}$.
  \item Find: overall efficiency $\eta$.
\end{itemize}

\textbf{First-law relation.}
\[
\eta = \frac{\dot W_{\text{out}}}{\dot Q_{\text{in}}}
\qquad\text{with}\qquad
\dot Q_{\text{in}} = \dot m_{\text{coal}}\,HV
\]

\textbf{Efficiency / COP / Carnot equation.} For a power plant, use efficiency $\eta$.

\textbf{Sequential working with units and time conversions.}
\[
60\,\text{tons/h} = 60{,}000\,\text{kg/h} = \frac{60{,}000}{3600}\,\text{kg/s} = 16.6667\,\text{kg/s}
\]
\[
\dot Q_{\text{in}} = (16.6667\,\text{kg/s})(30{,}000\,\text{kJ/kg}) = 500{,}000\,\text{kJ/s} = 500\,\text{MW}
\]
\[
\eta = \frac{150\,\text{MW}}{500\,\text{MW}} = 0.30
\]

\textbf{Boxed official answer.}
\[
\boxed{\eta = 0.30 = 30\%}
\]

\textbf{Feasibility / trap note.} The coal energy rate is much larger than the electrical output; 30\% is plausible for a steam plant. Do not confuse tons/h with kg/s.

\textbf{Exact source pages.} Question sheet p.~1; solution p.~1--2.

\subsection*{Q2}
\textbf{Rewritten question.} A heat pump used to heat a house has a COP of 2.5, meaning it delivers $2.5\,\text{kWh}$ of heat to the house for each $1\,\text{kWh}$ of electricity it consumes. Is this a violation of the first law of thermodynamics? Explain.

\textbf{Vague note.} The COP statement sounds larger than 1, but for a heat pump the delivered heat includes both the input work and heat absorbed from the outside air.

\textbf{Device/system type and reservoir boundary.} Closed cycle/steady cyclic device; hot reservoir is the house interior and cold reservoir is the outside air. The compressor work crosses the system boundary into the cycle.

\textbf{Assumptions.} Steady cyclic operation; negligible kinetic and potential energy changes; no mass crosses the boundary.

\textbf{Given / Find.}
\begin{itemize}
  \item Given: $\mathrm{COP}_{\mathrm{HP}}=2.5$.
  \item Find: whether the first law is violated.
\end{itemize}

\textbf{First-law relation.}
\[
\dot Q_H = \dot Q_L + \dot W_{\text{in}}
\]
For a heat pump,
\[
\mathrm{COP}_{\mathrm{HP}} = \frac{\dot Q_H}{\dot W_{\text{in}}}
\]

\textbf{Efficiency / COP / Carnot equation.} Heat pump COP is allowed to exceed 1 because it measures heat delivered per unit work input, not work output per heat input.

\textbf{Sequential working with units and time conversions.} Take $1\,\text{kWh}$ of electrical input.
\[
\dot W_{\text{in}} = 1\,\text{kWh}, \qquad \dot Q_H = 2.5\,\text{kWh}
\]
Then
\[
\dot Q_L = \dot Q_H - \dot W_{\text{in}} = 2.5 - 1.0 = 1.5\,\text{kWh}
\]
So the cycle moves $1.5\,\text{kWh}$ from outdoors into the house in addition to the $1\,\text{kWh}$ of work input.

\textbf{Boxed official answer.}
\[
\boxed{\text{No. The first law is not violated; the heat pump adds }1.5\,\text{kWh}\text{ from the outside air.}}
\]

\textbf{Feasibility / trap note.} The trap is assuming the delivered heat must equal the electrical input. It does not; the outside air supplies the rest.

\textbf{Exact source pages.} Question sheet p.~1; solution p.~3.

\subsection*{Q3}
\textbf{Rewritten question.} A refrigerator has a power input of $450\,\text{W}$ and a COP of 1.5. It is to cool five watermelons of mass $10\,\text{kg}$ each from $28^\circ\text{C}$ to $8^\circ\text{C}$. Treat the watermelons as water with $c=4.2\,\mathrm{kJ/(kg\,^{\circ}C)}$. How long will the cooling take, and is the result realistic or optimistic?

\textbf{Vague note.} The cooling load is the sensible heat removed from the melons, while the refrigerator extracts heat at a rate set by its COP.

\textbf{Device/system type and reservoir boundary.} Refrigerator cycle; cold reservoir is the watermelon load inside the cabinet, hot reservoir is the room air, and compressor work enters the cycle.

\textbf{Assumptions.} Watermelons are modeled as pure water; their temperature is uniform; refrigerator COP is constant; cabinet air thermal mass is negligible.

\textbf{Given / Find.}
\begin{itemize}
  \item Given: $\dot W_{\text{in}}=450\,\text{W}$, $\mathrm{COP}_R=1.5$, $m=5\times 10=50\,\text{kg}$, $T_1=28^\circ\text{C}$, $T_2=8^\circ\text{C}$, $c=4.2\,\mathrm{kJ/(kg\,^{\circ}C)}$.
  \item Find: time $t$.
\end{itemize}

\textbf{First-law relation.}
\[
\mathrm{COP}_R = \frac{\dot Q_L}{\dot W_{\text{in}}}
\qquad\Rightarrow\qquad
\dot Q_L = \mathrm{COP}_R\,\dot W_{\text{in}}
\]
and for the watermelon cooling load,
\[
Q = mc\Delta T
\]

\textbf{Efficiency / COP / Carnot equation.} For a refrigerator, use $\mathrm{COP}_R=\dot Q_L/\dot W_{\text{in}}$.

\textbf{Sequential working with units and time conversions.}
\[
\dot Q_L = 1.5\times 450\,\text{W} = 675\,\text{W} = 0.675\,\text{kJ/s}
\]
\[
Q = (50\,\text{kg})(4.2\,\mathrm{kJ/(kg\,^{\circ}C)})(28-8)^\circ\text{C}
= 4200\,\text{kJ}
\]
\[
t = \frac{Q}{\dot Q_L} = \frac{4200\,\text{kJ}}{0.675\,\text{kJ/s}} = 6222.22\,\text{s}
\]
\[
6222.22\,\text{s} = \frac{6222.22}{60}\,\text{min} = 103.70\,\text{min} \approx 1.73\,\text{h}
\]

\textbf{Boxed official answer.}
\[
\boxed{t \approx 6.22\times 10^3\,\text{s} \approx 104\,\text{min} \approx 1.73\,\text{h}}
\]

\textbf{Feasibility / trap note.} This is optimistic because real watermelons are not pure water and the cooling is not perfectly uniform; the actual time will usually be longer.

\textbf{Exact source pages.} Question sheet p.~1; solution p.~5--6.

\subsection*{Q4}
\textbf{Rewritten question.} Refrigerant-134a enters the condenser of a residential heat pump at $800\,\text{kPa}$ and $35^\circ\text{C}$ at a rate of $0.018\,\text{kg/s}$ and leaves at $800\,\text{kPa}$ as a saturated liquid. If the compressor consumes $1.2\,\text{kW}$ of power, determine (a) the COP of the heat pump and (b) the rate of heat absorption from the outside air.

\textbf{Vague note.} The condenser analysis gives the heat delivered to the house; then the cycle first law gives the outside-air heat absorption.

\textbf{Device/system type and reservoir boundary.} Residential heat pump cycle; condenser is the hot-side heat exchanger, evaporator is the cold-side heat exchanger, and the compressor work crosses the boundary into the cycle.

\textbf{Assumptions.} Steady flow through the condenser; negligible kinetic and potential energy changes; no shaft work in the condenser itself; property values are taken from R-134a tables.

\textbf{Given / Find.}
\begin{itemize}
  \item Given: $\dot m=0.018\,\text{kg/s}$, $P_1=800\,\text{kPa}$, $T_1=35^\circ\text{C}$, $P_2=800\,\text{kPa}$, saturated liquid at exit, $\dot W_{\text{in}}=1.2\,\text{kW}$.
  \item Find: (a) $\mathrm{COP}_{\mathrm{HP}}$, (b) $\dot Q_L$.
\end{itemize}

\textbf{First-law relation.}
For the condenser,
\[
\dot Q_H = \dot m\,(h_1-h_2)
\]
and for the full heat-pump cycle,
\[
\dot Q_H = \dot Q_L + \dot W_{\text{in}}
\]
so
\[
\mathrm{COP}_{\mathrm{HP}} = \frac{\dot Q_H}{\dot W_{\text{in}}},
\qquad
\dot Q_L = \dot Q_H - \dot W_{\text{in}}
\]

\textbf{Efficiency / COP / Carnot equation.} Heat pump COP is $\dot Q_H/\dot W_{\text{in}}$.

\textbf{Sequential working with units and time conversions.} From the R-134a tables at $800\,\text{kPa}$:
\[
h_1 \approx 272.7\,\text{kJ/kg}\quad\text{(slightly superheated at }35^\circ\text{C)}
\]
\[
h_2 \approx 95.5\,\text{kJ/kg}\quad\text{(saturated liquid)}
\]
Then
\[
\dot Q_H = (0.018\,\text{kg/s})(272.7-95.5)\,\text{kJ/kg}
= 3.19\,\text{kW}
\]
\[
\mathrm{COP}_{\mathrm{HP}} = \frac{3.19\,\text{kW}}{1.2\,\text{kW}} = 2.66
\]
\[
\dot Q_L = 3.19\,\text{kW} - 1.2\,\text{kW} = 1.99\,\text{kW}
\]

\textbf{Boxed official answer.}
\[
\boxed{\mathrm{COP}_{\mathrm{HP}} \approx 2.66, \qquad \dot Q_L \approx 1.99\,\text{kW}}
\]

\textbf{Feasibility / trap note.} The key trap is that the condenser outlet is saturated liquid, so the condenser enthalpy drop must be found from property tables; the COP is based on heat delivered, not compressor power alone.

\textbf{Exact source pages.} Question sheet p.~1; solution p.~7--13.

\subsection*{Q5}
\textbf{Rewritten question.} A Carnot heat engine operates between a source at $1000\,\text{K}$ and a sink at $300\,\text{K}$. If heat is supplied at a rate of $800\,\text{kJ/min}$, determine (a) the thermal efficiency and (b) the power output of the engine.

\textbf{Vague note.} Carnot means reversible, so the maximum possible efficiency is set only by the two reservoir temperatures.

\textbf{Device/system type and reservoir boundary.} Reversible heat engine; hot reservoir at $1000\,\text{K}$, cold reservoir at $300\,\text{K}$, and work leaves across the shaft boundary.

\textbf{Assumptions.} Carnot/reversible cycle; steady cyclic operation; negligible kinetic and potential energy changes.

\textbf{Given / Find.}
\begin{itemize}
  \item Given: $T_H=1000\,\text{K}$, $T_L=300\,\text{K}$, $\dot Q_H=800\,\text{kJ/min}$.
  \item Find: (a) $\eta_{\text{th}}$, (b) $\dot W_{\text{out}}$.
\end{itemize}

\textbf{First-law relation.}
\[
\eta_{\text{th}} = 1-\frac{\dot Q_L}{\dot Q_H} = \frac{\dot W_{\text{out}}}{\dot Q_H}
\]
For a Carnot engine,
\[
\eta_{\text{Carnot}} = 1-\frac{T_L}{T_H}
\]

\textbf{Efficiency / COP / Carnot equation.}
\[
\eta_{\text{Carnot}} = 1-\frac{300}{1000} = 0.70
\]

\textbf{Sequential working with units and time conversions.}
\[
800\,\text{kJ/min} = \frac{800}{60}\,\text{kJ/s} = 13.3333\,\text{kW}
\]
\[
\dot W_{\text{out}} = \eta_{\text{th}}\,\dot Q_H = 0.70\times 13.3333\,\text{kW} = 9.3333\,\text{kW}
\]

\textbf{Boxed official answer.}
\[
\boxed{\eta_{\text{th}}=70\%, \qquad \dot W_{\text{out}}\approx 9.33\,\text{kW}}
\]

\textbf{Feasibility / trap note.} The efficiency is not based on the heat input rate alone; it is limited by the hot and cold reservoir temperatures. The 70\% value is the Carnot maximum, not a guarantee for a real engine.

\textbf{Exact source pages.} Question sheet p.~2; solution p.~14--15.

\clearpage
\section{Final exam checks}
\begin{itemize}
\item The tutorial set contains 61 rewritten questions: Week 1 (10), Week 2 (10), Week 3 (14), Week 4 (12), Week 5 (10), Week 6 (5).
\item Quality $x$ is only valid inside a two-phase saturation region and must lie from 0 to 1.
\item A rigid tank has no moving-boundary work, but it may still receive heat or shaft/electrical work.
\item A filling or emptying tank is transient even if its inlet flow rate is constant.
\item A throttle is approximately isenthalpic, not isothermal or isentropic.
\item Carnot equations and ideal-gas equations require absolute temperature; ideal-gas pressure must be absolute.
\end{itemize}
\end{document}
